\documentclass[11pt]{article}
\usepackage[a4paper,margin=25mm]{geometry}
\usepackage[T1]{fontenc}
\usepackage[utf8]{inputenc}
\usepackage{lmodern,microtype}
\usepackage{amsmath,amssymb,amsthm,mathtools}
\usepackage{booktabs,longtable,tabularx,array,multirow}
\usepackage{enumitem}
\usepackage{xcolor}
\usepackage{hyperref}
\usepackage{fancyhdr}
\usepackage{listings}
\usepackage{url}
\usepackage{tikz}
\usetikzlibrary{arrows.meta,calc,angles,quotes}
\hypersetup{
  hypertexnames=false,
  colorlinks=true,
  linkcolor=blue!55!black,
  citecolor=blue!55!black,
  urlcolor=blue!55!black,
  pdftitle={Spherical Trigonometry and the Tetrahedral Cosine Law in Intuitionistic Primitive Recursive Arithmetic},
  pdfauthor={Research report prepared for Vladimir Reshetnikov},
  pdfsubject={Constructive spherical trigonometry in iPRA}
}
\pagestyle{fancy}\fancyhf{}
\fancyhead[L]{\small Spherical trigonometry in iPRA}
\fancyhead[R]{\small Constructive vector and finite-arithmetic proofs}
\fancyfoot[C]{\thepage}
\setlist{nosep}
\setlength{\parskip}{0.45em}
\setlength{\parindent}{0pt}
\allowdisplaybreaks
\sloppy

\newtheorem{theorem}{Theorem}[section]
\newtheorem{lemma}[theorem]{Lemma}
\newtheorem{proposition}[theorem]{Proposition}
\newtheorem{corollary}[theorem]{Corollary}
\theoremstyle{definition}
\newtheorem{definition}[theorem]{Definition}
\newtheorem{construction}[theorem]{Construction}
\newtheorem{remark}[theorem]{Remark}
\newtheorem{convention}[theorem]{Convention}

\newcommand{\N}{\mathbb N}
\newcommand{\Z}{\mathbb Z}
\newcommand{\Q}{\mathbb Q}
\newcommand{\R}{\mathbb R}
\newcommand{\C}{\mathbb C}
\newcommand{\ipra}{\mathsf{iPRA}}
\newcommand{\PRAAx}{\mathsf{PRAAx}}
\newcommand{\Prf}{\mathsf{Prf}}
\newcommand{\abs}[1]{\left|#1\right|}
\newcommand{\norm}[1]{\left\lVert#1\right\rVert}
\newcommand{\dotp}{\mathbin{\boldsymbol{\cdot}}}
\newcommand{\crossp}{\mathbin{\boldsymbol{\times}}}
\newcommand{\apart}{\mathrel{\#}}
\newcommand{\hav}{\operatorname{hav}}
\newcommand{\atanTwo}{\operatorname{atan2}}
\newcommand{\Area}{\operatorname{Area}}
\newcommand{\Vol}{\operatorname{Vol}}
\newcommand{\sgn}{\operatorname{sgn}}
\newcommand{\cis}{\operatorname{cis}}
\newcommand{\sech}{\operatorname{sech}}
\newcommand{\csch}{\operatorname{csch}}
\newcommand{\qturn}{q}
\newcommand{\code}[1]{\texttt{\detokenize{#1}}}
\newcommand{\cyc}{\operatorname{cyc}}
\newcommand{\BigO}{\mathcal O}
\DeclareMathOperator{\tr}{tr}
\DeclareMathOperator{\detm}{det}
\DeclareMathOperator{\arccot}{arccot}
\DeclareMathOperator{\arccsc}{arccsc}
\DeclareMathOperator{\arcsec}{arcsec}

\lstset{basicstyle=\ttfamily\small,breaklines=true,columns=fullflexible,frame=single}

\title{\textbf{Spherical Trigonometry and the Tetrahedral Cosine Law\\
       in Intuitionistic Primitive Recursive Arithmetic}\\[0.4em]
       \large Vector identities, constructive real names, a formula atlas, and a Lean completion route}
\author{Research report prepared for Vladimir Reshetnikov}
\date{15 November 2026}

\begin{document}
\maketitle

\begin{abstract}
We develop spherical trigonometry in intuitionistic primitive recursive
arithmetic (\(\ipra\)).  The core is deliberately finite.  A positively oriented
proper spherical triangle is represented by three unit vectors \(u,v,w\), and
all side and angle sines and cosines are recovered from dot products, cross
products and the scalar triple product.  Lagrange's identity and the \(3\times3\)
Gram determinant immediately yield the spherical cosine and sine laws.  The
polar triangle supplies the dual cosine law for angles and converts every
side formula into an angle formula and conversely.

From this kernel we prove the half-angle and half-side formulas, half-tangent
forms, haversine identities, Delambre--Gauss analogies, Napier's analogies,
the spherical law of tangents, cotangent and five-part rules, the complete
right and quadrantal triangle rules, Cagnoli's equation, differential
variation formulas, a quantitative planar limit, Girard's excess formula,
and the polygonal excess theorem.  For a tetrahedron, closure of the four
oriented face-area vectors gives the advertised cosine-law analogue for face
areas and dihedral angles; the facewise sine laws also give the tetrahedral
sine-product identity.

The phrase ``in \(\ipra\)'' is interpreted arithmetically rather than by silently
adding a completed real-number sort.  Inputs are fixed primitive-recursive
rapid dyadic names, supplied with regularity, magnitude, orientation and
apartness certificates.  Every division has a positive separation witness,
every square root is computed by bounded dyadic search, and each limit has an
explicit primitive-recursive modulus.  Most formulas reduce to checked
polynomial identities.  The bundled exact checker verifies 68 finite generator
identities in 14 families with zero residuals.  The archive contains the complete \LaTeX{} source,
this PDF, exact symbolic regression code, a formula-coverage ledger, and Lean
prototype files with a dependency-ordered path to a genuine
\(\Prf\;\PRAAx\) derivation.  The Lean files are an outline and algebraic
prototype, not a claim of completed object-level formalization.
\end{abstract}

\tableofcontents

\section{Scope, conventions, and the main theorem}
\label{sec:scope}

The formulas in the source pages are normally stated over ordinary real
numbers and use total-looking notation for quotients, inverse functions and
square roots.  Constructively these are not all statements of the same kind.
We therefore separate a finite algebraic kernel from the analytic presentation
layer.

\begin{definition}[Proof levels]
A formula in this report is one of the following.
\begin{enumerate}[label=\textup{(\roman*)}]
\item A \emph{polynomial identity}: after the displayed abbreviations are
expanded, it is an equality in a commutative ring.  Such a theorem needs no
order, excluded middle or real completeness.
\item A \emph{guarded rational identity}: all denominators carry explicit
apartness witnesses.  Its cross-multiplied form is polynomial.
\item A \emph{positive-root identity}: the radicand is proved nonnegative and
an interval or branch certificate selects the nonnegative square root.
\item A \emph{limit or differential identity}: every derivative, integral or
limit is represented by a fixed primitive-recursive rational-name program and
an explicit primitive-recursive modulus.
\end{enumerate}
\end{definition}

\begin{convention}[Proper oriented triangle]
Unless stated otherwise, \(A,B,C\in S^2\) are three pairwise non-antipodal
points joined by the shorter great-circle arcs.  Their side lengths satisfy
\[
 0<a,b,c<\pi,\qquad a+b+c<2\pi,
\]
and their interior angles satisfy \(0<A,B,C<\pi\).  We choose the orientation
so that the scalar triple product introduced in Section~\ref{sec:vector-model}
is positive.  In the arithmetic interpretation, each strict inequality is
positive data: a rational separation interval, not merely the negation of an
equality.
\end{convention}

The symbol \(q\) denotes a quarter turn and \(\pi=2q\).  The constructive
trigonometric packet cited in the bibliography constructs \(q\) as the unique
zero of cosine in \([1,2]\) by bounded sign search and proves
\(\sin q=1\), \(\cos q=0\).  Consequently no theorem below presupposes a
classical definition of \(\pi\).

\begin{theorem}[Constructive spherical-trigonometry atlas]
\label{thm:main-atlas}
For every fixed tuple of primitive-recursive regular dyadic names encoding a
proper oriented spherical triangle, together with fixed unit-vector,
orientation and apartness certificates, \(\ipra\) proves:
\begin{enumerate}[label=\textup{(\alph*)}]
\item the three spherical side cosine laws, the three angle cosine laws and the
spherical sine law;
\item the polar-triangle correspondences and involution;
\item all cyclic instances of the half-angle, half-side, half-tangent,
haversine, Delambre--Gauss, Napier, tangent, cotangent, five-part and Cagnoli
identities stated in this report;
\item the ten right-triangle and ten quadrantal-triangle rules, with explicit
guards on quotient forms;
\item the differential variation formulas for every fixed primitive-recursive
differentiable one-parameter family of proper triangles;
\item Girard's excess formula and the polygonal excess formula for every fixed
finite geodesic polygon with a certified triangulation;
\item the tetrahedral face-area cosine law and the tetrahedral sine-product
identity for every fixed nondegenerate tetrahedron with oriented face data.
\end{enumerate}
Every output proof can be generated from the finite input programs and their
certificates by a primitive-recursive metatheoretic proof transformer.
\end{theorem}

The quantification in Theorem~\ref{thm:main-atlas} is metatheoretic.  First-order
\(\ipra\) has numerical variables but neither a real-number sort nor function
variables.  Thus the theorem is a schema over fixed primitive-recursive terms;
it is not an assertion that \(\ipra\) contains a universal evaluator for all
primitive-recursive programs.  Section~\ref{sec:ipra} makes the translation
precise, and Appendix~\ref{app:ipra-target} displays a representative
object-language target.

\subsection{What is genuinely new work in the reduction}

The analytic sine/cosine layer can be reused from ProveIt's constructive
trigonometric-identities packet.  The new spherical contribution is the
following small kernel.
\begin{enumerate}
\item Encode a triangle by three unit vectors and retain the scalar triple
product as an orientation witness.
\item Derive all primary laws from two polynomial identities: Lagrange's
cross-product identity and the Gram determinant.
\item Use normalized face normals to construct the polar triangle without a
choice of coordinates.
\item Treat every familiar quotient formula as a cross-multiplied identity
plus explicit apartness data.
\item For tetrahedra, replace edge vectors by the four oriented face-area
vectors; their sum is zero, so the advertised cosine law is an ordinary norm
expansion.
\end{enumerate}
This organization makes the mathematical and formalization dependencies
transparent: the bulk of spherical trigonometry is finite algebra, not real
analysis.

\section{The arithmetic meaning of the geometry}
\label{sec:ipra}

\subsection{The base theory and the repository interface}

We use intuitionistic primitive recursive arithmetic with symbols generated by
zero, successor, projections, composition and primitive recursion, defining
equations for those symbols, equality axioms, successor separation and
injectivity, and induction for quantifier-free formulas.  In ProveIt's
\code{ConciseIPRA} development these ingredients are represented by the types
\code{PR k}, \code{PTm n}, \code{Fm}, the proof relation \code{Prf}, and the
theory \code{PRAAx}.  The derived rule \code{PRA.qfInd} requires an explicit
proof that its induction formula is quantifier free.  The files
\code{BiInterpretation.lean} and \code{NativeBiInterpretation.lean} provide
transport only after an actual object-language derivation has been built; an
ambient theorem over Lean's \code{Real} type does not automatically cross this
boundary.

All numerical data below are ultimately natural-number tuples.  Signed
integers may be coded as pairs \((p,n)\), denoting \(p-n\).  A rational is a
signed numerator together with a positive denominator.  Unreduced fractions
are sufficient: equality and order are tested by cross multiplication, and
all loops used below are bounded.

\subsection{Rapid dyadic names}

\begin{definition}[Regular rapid name]
A rapid dyadic name is a fixed primitive-recursive function
\(x:\N\to\Q\), together with a proof in \(\ipra\) that
\begin{equation}
  \abs{x(n+1)-x(n)}\le 2^{-n-4}                         \tag{2.1}
\end{equation}
for every \(n\), and an integer \(B_x\ge1\) satisfying
\(\abs{x(n)}\le B_x\).  The associated interval is
\begin{equation}
 I_x(n)=\bigl[x(n)-2^{-n-3},\,x(n)+2^{-n-3}\bigr].       \tag{2.2}
\end{equation}
\end{definition}

A finite geometric-sum induction gives, for \(m\ge n\),
\[
 \abs{x(m)-x(n)}<2^{-n-3}.
\]
Thus every later value lies in \(I_x(n)\).  Equality of names is the numerical
formula
\begin{equation}
 x=_\R y\quad:\Longleftrightarrow\quad
 \forall n\;\abs{x(n+6)-y(n+6)}\le2^{-n}.               \tag{2.3}
\end{equation}
The constants in (2.1)--(2.3) are deliberately generous.  Their purpose is to
make every precision shift primitive recursive and easy to audit.

\begin{definition}[Apartness and strict order]
An apartness witness \(x\apart0\) consists of natural numbers \(n,k\) and a
sign \(\sigma\in\{-1,+1\}\) such that the entire interval \(I_x(n)\) lies in
\([2^{-k},\infty)\) when \(\sigma=+1\), or in
\(({-}\infty,-2^{-k}]\) when \(\sigma=-1\).  A witness \(x<y\) is a finite
precision at which the upper endpoint of \(I_x\) is strictly below the lower
endpoint of \(I_y\).
\end{definition}

Apartness is stronger than \(x\ne0\) and is exactly the information needed to
compute a reciprocal.  This distinction is essential.  No proof below infers
a denominator bound from a bare negated equality.

\subsection{Arithmetic closure and bounded roots}

Addition and negation use fixed precision shifts.  On bounded inputs,
multiplication follows from
\begin{equation}
 \abs{xy-x'y'}\le \abs{x}\abs{y-y'}+\abs{y'}\abs{x-x'}.
                                                               \tag{2.4}
\end{equation}
If \(\abs{x}\ge d>0\) is supplied, reciprocal stability is
\begin{equation}
 \abs{x^{-1}-y^{-1}}
 =\frac{\abs{x-y}}{\abs{x}\abs{y}}
 \le \frac{2\abs{x-y}}{d^2}                                 \tag{2.5}
\end{equation}
whenever the input error is at most \(d/2\).  The right-hand side determines a
primitive-recursive precision transformer.

For a nonnegative rational \(r=p/q\), requested precision \(k\), and
\(S=2^k\), scan the bounded interval
\(0\le j\le(p+1)S\) and retain the last \(j\) with
\begin{equation}
 qj^2\le pS^2.                                                \tag{2.6}
\end{equation}
The endpoint fails (2.6), and the quantifier-free loop invariant says that the
retained value succeeds while its successor fails.  Hence
\[
 \frac jS\le\sqrt{p/q}<\frac{j+1}{S}.
\]
This is the square-root routine used for vector norms.  For nonnegative real
names, apply it to a sufficiently accurate rational enclosure.  The estimate
\begin{equation}
 \abs{\sqrt{x}-\sqrt{y}}\le\sqrt{\abs{x-y}}                 \tag{2.7}
\end{equation}
shows that doubling the requested input precision suffices near zero.  No
unbounded search and no exact real comparison occurs.

All finite loops just described have decidable rational guards.  Their
invariants can be expressed by primitive-recursive characteristic equations,
so quantifier-free induction applies.  Universal quantification over the final
precision parameter is introduced only after the finite estimate has been
proved.

\subsection{Constructive trigonometric names}
\label{subsec:trig-names}

The ProveIt trigonometric-identities packet supplies an analytic realization
that can be reused verbatim.  For a Gaussian rational \(z\) with
\(\norm{z}_1\le B\), use the finite exponential sum
\[
 E_N(z)=\sum_{j=0}^{N-1}\frac{z^j}{j!},
 \qquad N(B,n)=2B^2+2B+n+4.                                 \tag{2.8}
\]
Beyond index \(2B\), the majorant terms have ratio at most \(1/2\), and a
finite geometric estimate gives a tail below \(2^{-n-3}\).  A finite
Cauchy-product identity proves
\(\exp(z+w)=\exp z\exp w\).  Then
\begin{equation}
 \sin x=\frac{\exp(ix)-\exp(-ix)}{2i},\qquad
 \cos x=\frac{\exp(ix)+\exp(-ix)}2                         \tag{2.9}
\end{equation}
are primitive-recursive rapid names and satisfy the addition formulas and
\(\sin^2x+\cos^2x=1\).

A quarter turn \(q\in[1,2]\) is constructed by a bounded rational bisection
using explicit alternating-series signs for \(\cos1\) and \(\cos2\) and a
quantitative monotonicity estimate.  It satisfies
\begin{equation}
 \cos q=0,\qquad \sin q=1,
 \qquad \pi:=2q.                                             \tag{2.10}
\end{equation}
Thus every occurrence of \(\pi\), periodicity, or a supplementary angle in this
report has a concrete arithmetic realization.

\subsection{Angles from circle-valued data}

Most proofs need only a circle-valued angle pair \((c,s)\) satisfying
\(c^2+s^2=1\), with a sign witness for \(s\).  When an actual angle name is
desired, use two rationally separated charts.  Given a unit pair with
\(s>0\), compute once a rational approximation \(z\) to \(c\) with error at
most \(1/4\).
\begin{itemize}
\item If \(z\ge0\), then \(1+c\ge1/2\) and define
\[
 \theta=2\arctan\!\left(\frac{s}{1+c}\right).
\]
\item If \(z<0\), then \(1-c\ge1/2\) and define
\[
 \theta=\pi-2\arctan\!\left(\frac{s}{1-c}\right).
\]
\end{itemize}
The chart decision is finite and retained at every output precision.  The
arctangent is evaluated by range reduction and an alternating rational series
with an explicit primitive-recursive tail bound.  The Cayley identities
\[
 \cos(2\arctan t)=\frac{1-t^2}{1+t^2},\qquad
 \sin(2\arctan t)=\frac{2t}{1+t^2}                           \tag{2.11}
\]
prove that the resulting name realizes \((c,s)\).  For side lengths and
interior angles in \((0,\pi)\), the positive sine witness fixes the branch
uniquely.  This is a constructive \(\atanTwo(s,c)\) operation, not a global
continuous choice on the whole circle.

\subsection{The schema, not a universal evaluator}

For a fixed spherical triangle, the input consists of finitely many particular
\code{PR} syntax trees for coordinate names, bounds and certificates.  The
output operations are new particular \code{PR} trees obtained by composition,
primitive recursion and bounded searches.  The metatheory then constructs a
finite derivation
\[
 \Prf\;\PRAAx\;[]\;(\text{hypotheses}\to\text{conclusion}). \tag{2.12}
\]
There is no quantification over an arbitrary program code followed by an
unbounded universal evaluation operation.  Such an evaluator cannot enumerate
all primitive-recursive functions.  This fixed-program schema is the precise
sense in which the theorems are proved in first-order \(\ipra\).

\begin{remark}[Ambient Lean versus object-level arithmetic]
A proof by \code{ring} over a Lean commutative ring is valuable: it checks the
finite identity that the arithmetic derivation must replay.  It is not itself
a term of \(\Prf\;\PRAAx\).  The final step must either write the corresponding
equational derivation directly or invoke a polynomial normalizer whose
soundness has itself been proved inside the object theory.  The Lean roadmap in
Section~\ref{sec:formalization} preserves this distinction explicitly.
\end{remark}

\section{The finite vector model of a spherical triangle}
\label{sec:vector-model}

Let \(u,v,w\in\R^3\) be the unit vectors pointing to the vertices
\(A,B,C\), respectively.  The side \(a\) is opposite \(A\), hence it is the
shorter central angle between \(v\) and \(w\); cyclically for \(b,c\).  Put
\begin{equation}
 x=v\dotp w=\cos a,\qquad
 y=w\dotp u=\cos b,\qquad
 z=u\dotp v=\cos c.                                         \tag{3.1}
\end{equation}
Define the positive side sines
\begin{equation}
 s_a=\norm{v\crossp w},\qquad
 s_b=\norm{w\crossp u},\qquad
 s_c=\norm{u\crossp v}.                                    \tag{3.2}
\end{equation}
Since the sides lie strictly between \(0\) and \(\pi\), these norms are
positively apart from zero.  Lagrange's identity gives
\begin{equation}
 s_a^2=1-x^2,\qquad s_b^2=1-y^2,\qquad s_c^2=1-z^2.          \tag{3.3}
\end{equation}
Finally put
\begin{equation}
 \tau=\detm[u\ v\ w]=u\dotp(v\crossp w)>0.                 \tag{3.4}
\end{equation}
The sign of \(\tau\) is the orientation certificate.  Nondegeneracy supplies a
rational lower bound for it.

\subsection{The Gram identity}

\begin{lemma}[Gram determinant]
\label{lem:gram}
For unit \(u,v,w\),
\begin{equation}
 \tau^2=1+2xyz-x^2-y^2-z^2.                                 \tag{3.5}
\end{equation}
\end{lemma}

\begin{proof}
The square of a scalar triple product is the determinant of the Gram matrix:
\[
 \tau^2=
 \detm\!\begin{pmatrix}
 1&z&y\\ z&1&x\\ y&x&1
 \end{pmatrix}.
\]
Expanding this \(3\times3\) determinant gives (3.5).  In coordinates the claim
is a polynomial identity with integer coefficients, so its arithmetic proof
needs only distributivity and commutativity.
\end{proof}

A useful factorization is
\begin{equation}
 \tau^2=(1-y^2)(1-z^2)-(x-yz)^2
       =s_b^2s_c^2-(x-yz)^2,                               \tag{3.6}
\end{equation}
with two cyclic companions.

\subsection{Tangent directions and angle pairs}

At the vertex \(A\), the unit tangent direction toward \(B\) is
\begin{equation}
 t_{AB}=\frac{v-zu}{s_c},                                   \tag{3.7}
\end{equation}
and the unit tangent direction toward \(C\) is
\begin{equation}
 t_{AC}=\frac{w-yu}{s_b}.                                   \tag{3.8}
\end{equation}
Indeed, the numerators are orthogonal to \(u\), and their squared norms are
\(1-z^2=s_c^2\) and \(1-y^2=s_b^2\).  Define the trigonometric pair of the
interior angle \(A\) by
\begin{align}
 c_A&=t_{AB}\dotp t_{AC}
     =\frac{x-yz}{s_bs_c},                                  \tag{3.9}\\
 s_A&=u\dotp(t_{AB}\crossp t_{AC})
     =\frac{\tau}{s_bs_c}.                                  \tag{3.10}
\end{align}
Cyclically,
\begin{align}
 c_B&=\frac{y-zx}{s_cs_a},&s_B&=\frac{\tau}{s_cs_a},        \tag{3.11}\\
 c_C&=\frac{z-xy}{s_as_b},&s_C&=\frac{\tau}{s_as_b}.        \tag{3.12}
\end{align}

\begin{lemma}[Angle pairs are unit pairs]
\label{lem:angle-unit}
For each vertex,
\[
 c_A^2+s_A^2=1,
 \qquad c_B^2+s_B^2=1,
 \qquad c_C^2+s_C^2=1.
\]
Moreover \(s_A,s_B,s_C>0\).
\end{lemma}

\begin{proof}
Equation (3.6) says
\((x-yz)^2+\tau^2=s_b^2s_c^2\).  Divide by the positively separated
quantity \(s_b^2s_c^2\).  The other equations are cyclic.  Positivity follows
from \(\tau>0\) and the positive side sines.
\end{proof}

The constructive angle-chart procedure of Section~\ref{subsec:trig-names}
therefore produces unique angle names \(A,B,C\in(0,\pi)\) with
\(\cos A=c_A\), \(\sin A=s_A\), and cyclically.  Notice that the algebraic
proofs can work entirely with the pairs \((c_A,s_A)\); actual radian-valued
angle names are needed only for formulas involving sums, halves or area.

\subsection{A compact certificate for a spherical triangle}

For object-level formalization it is convenient to regard the following as the
primitive certificate:
\begin{equation}
 \mathcal T=(u,v,w;x,y,z;s_a,s_b,s_c;\tau;
             c_A,s_A,c_B,s_B,c_C,s_C),                     \tag{3.13}
\end{equation}
with the equations (3.1)--(3.12), unit-vector equations, and positive rational
lower bounds for \(s_a,s_b,s_c,\tau\).  Every equation is polynomial after
clearing the displayed positive denominators.  A checker for \(\mathcal T\)
uses only signed rational arithmetic, bounded square-root brackets and finite
comparisons.  It does not decide equality of arbitrary real names.

\begin{proposition}[Finite-coordinate sufficiency]
Every division-free spherical identity proved through
Section~\ref{sec:cotangent-special} follows from the polynomial equations in
\(\mathcal T\), the addition laws for a constructive trigonometric system, and
finitely many positive denominator witnesses.
\end{proposition}

\begin{proof}
The primary laws are immediate consequences of (3.9)--(3.12).  The remaining
identities are obtained by finite addition, multiplication, substitution and
cross multiplication.  Half formulas additionally use the positive branches
fixed by \(0<a,b,c,A,B,C<\pi\).  No compactness, choice, maximum principle or
uncountable construction enters.
\end{proof}

\section{The spherical cosine and sine laws}
\label{sec:core-laws}

\subsection{The side cosine laws}

\begin{theorem}[Spherical law of cosines for sides]
\label{thm:side-cosine}
For a proper spherical triangle,
\begin{align}
 \cos a&=\cos b\cos c+\sin b\sin c\cos A,                 \tag{4.1}\\
 \cos b&=\cos c\cos a+\sin c\sin a\cos B,                 \tag{4.2}\\
 \cos c&=\cos a\cos b+\sin a\sin b\cos C.                 \tag{4.3}
\end{align}
Equivalently, without division,
\begin{align}
 s_bs_c c_A&=x-yz,&
 s_cs_a c_B&=y-zx,&
 s_as_b c_C&=z-xy.                                        \tag{4.4}
\end{align}
\end{theorem}

\begin{proof}
At \(A\), use the tangent vectors (3.7)--(3.8):
\[
 \cos A=t_{AB}\dotp t_{AC}
 =\frac{(v-zu)\dotp(w-yu)}{s_cs_b}
 =\frac{x-yz}{s_bs_c}.
\]
Multiplication by \(s_bs_c\) gives
\(x=yz+s_bs_c\cos A\), which is (4.1).  Cyclic relabeling gives
(4.2)--(4.3).  Every step before the final presentation is a polynomial dot-
product expansion; division uses the supplied positive side-sine witnesses.
\end{proof}

For a sphere of radius \(R>0\), the dimensionless central angles are
\(a/R,b/R,c/R\).  Thus (4.1) becomes
\[
 \cos\frac aR=\cos\frac bR\cos\frac cR+
 \sin\frac bR\sin\frac cR\cos A.                           \tag{4.5}
\]
The proof is unchanged after normalizing the radius vector by \(R\).

\begin{corollary}[Rearranged cosine formulas]
\label{cor:rearranged-cosine}
With the usual positive denominator witnesses,
\begin{align}
 \cos A&=\frac{\cos a-\cos b\cos c}{\sin b\sin c},         \tag{4.6}\\
 \cos a&=\frac{\cos A+\cos B\cos C}{\sin B\sin C}.         \tag{4.7}
\end{align}
The remaining formulas are cyclic.
\end{corollary}

The second formula is proved below as the dual cosine law.  Constructively,
(4.6) is shorthand for (4.4) plus the apartness of \(\sin b\sin c\).

\subsection{The spherical sine law}

\begin{theorem}[Spherical law of sines]
\label{thm:sine-law}
For a proper spherical triangle,
\begin{equation}
 \frac{\sin A}{\sin a}
 =\frac{\sin B}{\sin b}
 =\frac{\sin C}{\sin c}
 =\frac{\tau}{\sin a\sin b\sin c}.                         \tag{4.8}
\end{equation}
Its division-free content is
\begin{equation}
 \sin A\sin b=\sin B\sin a,
 \quad
 \sin B\sin c=\sin C\sin b,
 \quad
 \sin C\sin a=\sin A\sin c.                              \tag{4.9}
\end{equation}
\end{theorem}

\begin{proof}
Equation (3.10) gives
\[
 \sin A=\frac{\tau}{s_bs_c}.
\]
Dividing by \(s_a=\sin a\) yields the last expression in (4.8).
The two cyclic formulas have the same numerator and denominator.  For the
strictly algebraic version, cross multiply; for example
\[
 \sin A\sin b=\frac{\tau}{s_bs_c}s_b=\frac\tau{s_c}
 =\frac{\tau}{s_cs_a}s_a=\sin B\sin a.
\]
All canceled quantities have positive apartness witnesses.
\end{proof}

\begin{corollary}[Scalar-triple-product and volume forms]
\label{cor:sine-volume}
If \(V\) is the Euclidean volume of the tetrahedron with vertices
\(0,u,v,w\), then \(6V=\tau\), and
\begin{equation}
 \frac{\sin A}{\sin a}
 =\frac{6V}{\sin a\sin b\sin c}.                           \tag{4.10}
\end{equation}
This is the vector/triple-product form commonly used to prove the spherical
sine law.
\end{corollary}

\subsection{The angle cosine laws}

\begin{theorem}[Supplemental spherical law of cosines]
\label{thm:angle-cosine}
The angles and opposite sides satisfy
\begin{align}
 \cos A&=-\cos B\cos C+\sin B\sin C\cos a,                \tag{4.11}\\
 \cos B&=-\cos C\cos A+\sin C\sin A\cos b,                \tag{4.12}\\
 \cos C&=-\cos A\cos B+\sin A\sin B\cos c.                \tag{4.13}
\end{align}
Equivalently,
\begin{equation}
 \cos a=\frac{\cos A+\cos B\cos C}{\sin B\sin C},        \tag{4.14}
\end{equation}
with cyclic companions and explicit denominator witnesses.
\end{theorem}

\begin{proof}
We give a direct finite-algebra proof; Section~\ref{sec:polar} gives the more
conceptual proof by polar duality.  Substitute (3.11)--(3.12) into the
right-hand side of (4.11).  After multiplication by the positive quantity
\(s_a^2s_bs_c\), the desired equality becomes
\begin{equation}
 -(y-zx)(z-xy)+x\tau^2=(1-x^2)(x-yz).                      \tag{4.15}
\end{equation}
Insert the Gram identity (3.5); expansion reduces both sides to the same
integer polynomial.  Dividing by \(s_a^2s_bs_c\) gives (4.11).  The other
identities are cyclic.  Solving (4.11) for \(\cos a\) gives (4.14).
\end{proof}

\subsection{The complete core as a finite ideal}

The primary spherical theory can therefore be generated by the polynomial
relations
\begin{align*}
&s_a^2+x^2=1,\quad s_b^2+y^2=1,\quad s_c^2+z^2=1,\\
&\tau^2=1+2xyz-x^2-y^2-z^2,\\
&s_bs_c c_A=x-yz,\quad s_cs_a c_B=y-zx,\quad
 s_as_b c_C=z-xy,\\
&s_bs_c s_A=\tau,\quad s_cs_a s_B=\tau,\quad
 s_as_b s_C=\tau.                                          \tag{4.16}
\end{align*}
Every formula in this section is a consequence of (4.16) by polynomial
normalization and guarded cancellation.  This finite presentation is the core
object to formalize in Lean and in \(\ipra\).

\section{The polar triangle and side--angle duality}
\label{sec:polar}

The polar triangle is the mechanism behind the perfect duality of spherical
trigonometry.  In the present vector model it is constructive and explicit.

\begin{definition}[Positive polar triangle]
Define
\begin{equation}
 u^*=\frac{v\crossp w}{s_a},\qquad
 v^*=\frac{w\crossp u}{s_b},\qquad
 w^*=\frac{u\crossp v}{s_c}.                               \tag{5.1}
\end{equation}
The signs are fixed by
\(u\dotp u^*=v\dotp v^*=w\dotp w^*=\tau/(s_a,s_b,s_c)>0\),
so no binary choice of antipodes remains.
\end{definition}

Each polar vertex is a unit vector and is perpendicular to the great circle
containing the opposite original side.

\begin{theorem}[Polar side and angle correspondences]
\label{thm:polar-correspondence}
Let \(a^*,b^*,c^*\) and \(A^*,B^*,C^*\) denote the sides and angles of the
positive polar triangle.  Then
\begin{equation}
 a^*=\pi-A,
 \qquad b^*=\pi-B,
 \qquad c^*=\pi-C,                                         \tag{5.2}
\end{equation}
and
\begin{equation}
 A^*=\pi-a,
 \qquad B^*=\pi-b,
 \qquad C^*=\pi-c.                                         \tag{5.3}
\end{equation}
The polar operation is involutive: \((u^*)^*=u\), cyclically.
\end{theorem}

\begin{proof}
By Lagrange's four-vector identity,
\begin{align*}
 v^*\dotp w^*
 &=\frac{(w\crossp u)\dotp(u\crossp v)}{s_bs_c}\\
 &=\frac{(w\dotp u)(u\dotp v)-(w\dotp v)(u\dotp u)}{s_bs_c}\\
 &=\frac{yz-x}{s_bs_c}=-\cos A=\cos(\pi-A).
\end{align*}
The side \(a^*\) lies in \((0,\pi)\), and its sine is positive, so the branch is
uniquely \(a^*=\pi-A\).  The other side identities are cyclic.

For involution, use the vector triple-product identity:
\begin{equation}
 (w\crossp u)\crossp(u\crossp v)
 =u\bigl((w\crossp u)\dotp v\bigr)=\tau u.                 \tag{5.4}
\end{equation}
After division by \(s_bs_c\), the norm of the left-hand side is
\(\tau/(s_bs_c)=s_A=\sin a^*\).  Normalizing therefore returns \(u\).
The angle correspondence (5.3) follows by applying (5.2) to the polar triangle
and using involution.
\end{proof}

All steps in this proof are polynomial before positive normalization.  In
particular, polar duality does not require a compactness argument or a choice
of one of two normals: the orientation witness \(\tau>0\) fixes the sign.

\begin{corollary}[Duality principle]
\label{cor:duality-principle}
Any identity valid for all proper spherical triangles remains valid after the
simultaneous substitution
\begin{equation}
 (a,b,c;A,B,C)
 \longmapsto
 (\pi-A,\pi-B,\pi-C;\pi-a,\pi-b,\pi-c),                   \tag{5.5}
\end{equation}
with the corresponding changes
\(\sin(\pi-x)=\sin x\) and \(\cos(\pi-x)=-\cos x\).
\end{corollary}

Applying the side cosine law to the polar triangle immediately gives
\[
 -\cos A=\cos B\cos C-\sin B\sin C\cos a,
\]
which is Theorem~\ref{thm:angle-cosine}.  Thus the supplemental cosine law is
not an independent analytic theorem; it is the ordinary cosine law under the
finite polar construction.

\subsection{Half-sums under polarity}

Put
\begin{equation}
 p=\frac{a+b+c}{2},\qquad P=\frac{A+B+C}{2}.                \tag{5.6}
\end{equation}
For the polar triangle,
\begin{equation}
 p^*=\frac{a^*+b^*+c^*}{2}=\frac{3\pi}{2}-P,
 \qquad
 P^*=\frac{3\pi}{2}-p.                                    \tag{5.7}
\end{equation}
Consequently
\begin{align}
 \sin(p^*-b^*)&=\cos(P-B),&
 \sin p^*&=-\cos P,                                       \tag{5.8}\\
 \sin b^*&=\sin B,&
 \cos(P^*-B^*)&=\sin(p-b),                                \tag{5.9}
\end{align}
with cyclic companions.  These elementary substitutions turn the angle
half-formulas into the side half-formulas in Section~\ref{sec:half-formulas}.

\begin{corollary}[Polar excess and perimeter]
If \(E=A+B+C-\pi\) is the original spherical excess and \(E^*\) the polar
excess, then
\begin{equation}
 E^*=2\pi-(a+b+c).                                         \tag{5.10}
\end{equation}
Thus the proper-triangle condition \(a+b+c<2\pi\) is equivalent to positive
excess of the polar triangle.
\end{corollary}

\section{Half-angle, half-side, half-tangent, and haversine formulas}
\label{sec:half-formulas}

Recall
\[
 p=\frac{a+b+c}{2},\qquad P=\frac{A+B+C}{2}.
\]
Every side and angle half lies in \((0,\pi/2)\).  Consequently all square roots
in this section mean the nonnegative root and no independent sign choice is
hidden in the notation.

\subsection{Half-angle formulas}

\begin{theorem}[Spherical half-angle formulas]
\label{thm:half-angle}
For the angle \(A\),
\begin{align}
 \sin^2\frac A2
 &=\frac{\sin(p-b)\sin(p-c)}{\sin b\sin c},                \tag{6.1}\\
 \cos^2\frac A2
 &=\frac{\sin p\sin(p-a)}{\sin b\sin c},                  \tag{6.2}\\
 \tan^2\frac A2
 &=\frac{\sin(p-b)\sin(p-c)}{\sin p\sin(p-a)}.            \tag{6.3}
\end{align}
Hence
\begin{align}
 \sin\frac A2
 &=\sqrt{\frac{\sin(p-b)\sin(p-c)}{\sin b\sin c}},        \tag{6.4}\\
 \cos\frac A2
 &=\sqrt{\frac{\sin p\sin(p-a)}{\sin b\sin c}},           \tag{6.5}\\
 \tan\frac A2
 &=\sqrt{\frac{\sin(p-b)\sin(p-c)}{\sin p\sin(p-a)}}.     \tag{6.6}
\end{align}
The formulas for \(B,C\) are cyclic.
\end{theorem}

\begin{proof}
Subtract the side cosine law from the cosine difference formula:
\begin{align*}
 \cos(b-c)-\cos a
 &=\sin b\sin c(1-\cos A)\\
 &=2\sin b\sin c\sin^2(A/2).
\end{align*}
On the other hand, product-to-sum gives
\[
 \cos(b-c)-\cos a=2\sin(p-b)\sin(p-c).
\]
This proves (6.1).  Similarly,
\begin{align*}
 \cos a-\cos(b+c)
 &=\sin b\sin c(1+\cos A)\\
 &=2\sin b\sin c\cos^2(A/2),
\end{align*}
whereas
\[
 \cos a-\cos(b+c)=2\sin p\sin(p-a).
\]
This proves (6.2).  Divide (6.1) by (6.2), using the positive apartness of
\(\cos(A/2)\), to obtain (6.3).  Positive-root uniqueness gives
(6.4)--(6.6).  The proof uses only the cosine law and finite trigonometric
identities already available in \(\ipra\).
\end{proof}

A useful common-factor form is obtained by setting
\begin{equation}
 k=\frac{1}{\sqrt{\sin p\sin(p-a)\sin(p-b)\sin(p-c)}}.     \tag{6.7}
\end{equation}
Then
\begin{equation}
 \tan\frac A2=k\sin(p-b)\sin(p-c),                         \tag{6.8}
\end{equation}
with cyclic companions.

\subsection{Half-side formulas}

\begin{theorem}[Spherical half-side formulas]
\label{thm:half-side}
For the side \(a\),
\begin{align}
 \sin^2\frac a2
 &=-\frac{\cos P\cos(P-A)}{\sin B\sin C},                 \tag{6.9}\\
 \cos^2\frac a2
 &=\frac{\cos(P-B)\cos(P-C)}{\sin B\sin C},                \tag{6.10}\\
 \tan^2\frac a2
 &=-\frac{\cos P\cos(P-A)}{\cos(P-B)\cos(P-C)}.           \tag{6.11}
\end{align}
Consequently
\begin{equation}
 \tan\frac a2=K\cos(P-A),                                 \tag{6.12}
\end{equation}
where
\begin{equation}
 K=\sqrt{-\frac{\cos P}
 {\cos(P-A)\cos(P-B)\cos(P-C)}}.                          \tag{6.13}
\end{equation}
The formulas for \(b,c\) are cyclic.
\end{theorem}

\begin{proof}
Apply Theorem~\ref{thm:half-angle} to the polar triangle.  Its angle at
\(A^*\) is \(\pi-a\), its opposite side is \(\pi-A\), and its half-sums are
given by (5.7).  Thus
\[
 \cos^2\frac{A^*}{2}=\sin^2\frac a2
 =\frac{\sin p^*\sin(p^*-a^*)}{\sin b^*\sin c^*}
 =-\frac{\cos P\cos(P-A)}{\sin B\sin C},
\]
which is (6.9).  Likewise
\[
 \sin^2\frac{A^*}{2}=\cos^2\frac a2
 =\frac{\cos(P-B)\cos(P-C)}{\sin B\sin C},
\]
giving (6.10).  Their quotient gives (6.11).  Equation (6.12) is simply the
positive square root of (6.11) after extracting the common factor (6.13).
Every sign is certified by the proper-triangle hypotheses and the polar
construction.
\end{proof}

\begin{remark}[The Wikipedia half-side formula]
Equation (6.11), followed by the positive square root, is the standard
half-side formula
\[
 \tan\frac a2=
 \sqrt{-\frac{\cos P\cos(P-A)}
 {\cos(P-B)\cos(P-C)}}.
\]
The constructive statement includes the proof that the radicand is
nonnegative and a branch certificate selecting \(0<a/2<\pi/2\).
\end{remark}

\subsection{Direct vector half-angle formulas}

The vector certificate yields particularly short expressions:
\begin{align}
 \tan\frac A2
 &=\frac{\sin A}{1+\cos A}
 =\frac{\tau}{s_bs_c+x-yz},                                \tag{6.14}\\
 \tan\frac a2
 &=\frac{\sin a}{1+\cos a}
 =\frac{s_a}{1+x}.                                         \tag{6.15}
\end{align}
The denominators are positive because \(A,a\in(0,\pi)\).  These formulas are
well suited to computation: they use one square root for each side sine and no
inverse trigonometric function.

\subsection{Half-tangent rational forms}

Put
\[
 t_a=\tan\frac a2,
 \quad t_b=\tan\frac b2,
 \quad t_c=\tan\frac c2,
 \qquad
 t_A=\tan\frac A2,
 \quad t_B=\tan\frac B2,
 \quad t_C=\tan\frac C2.
\]
Tangent addition, together with Theorems~\ref{thm:half-angle} and
\ref{thm:half-side}, gives the entirely rational forms
\begin{equation}
 t_A^2=
 \frac{(t_a-t_b+t_c+t_at_bt_c)(t_a+t_b-t_c+t_at_bt_c)}
 {(t_a+t_b+t_c-t_at_bt_c)(-t_a+t_b+t_c+t_at_bt_c)},        \tag{6.16}
\end{equation}
and
\begin{equation}
 t_a^2=
 \frac{(t_Bt_C+t_Ct_A+t_At_B-1)
       (-t_Bt_C+t_Ct_A+t_At_B+1)}
 {(t_Bt_C-t_Ct_A+t_At_B+1)
       (t_Bt_C+t_Ct_A-t_At_B+1)}.                          \tag{6.17}
\end{equation}
All displayed denominator factors carry apartness witnesses.  The remaining
five equations are obtained cyclically.  To verify (6.16), write every
\(\sin(p-a)\) and \(\sin p\) using the tangent-of-a-sum formula and cancel the
common positive products \(\prod(1+t_x^2)^{-1/2}\).  Formula (6.17) is its polar
dual.  This derivation is finite rational algebra, so after cross multiplication
it is directly suitable for an exact certificate checker.

\subsection{Haversine formulas}

Define
\begin{equation}
 \hav t=\sin^2\frac t2=\frac{1-\cos t}{2}.                 \tag{6.18}
\end{equation}
The side cosine law becomes the numerically stable navigation formula
\begin{equation}
 \hav a=\hav(b-c)+\sin b\sin c\,\hav A,                    \tag{6.19}
\end{equation}
with cyclic companions.  Indeed,
\[
 \cos(b-c)-\cos a=\sin b\sin c(1-\cos A),
\]
and division by \(2\) gives (6.19).  The dual angle form is
\begin{equation}
 \hav A=\hav(\pi-B-C)+\sin B\sin C\,\hav a.                \tag{6.20}
\end{equation}
When \(\sin b\sin c\apart0\), one may solve for the angle:
\begin{equation}
 \hav A=\frac{\hav a-\hav(b-c)}{\sin b\sin c}.             \tag{6.21}
\end{equation}
Equations (6.18)--(6.21), their cyclic forms, and their polar duals account for
the standard haversine family without any loss of information from the cosine
law.

\section{Delambre--Gauss analogies, Napier analogies, and the tangent law}
\label{sec:analogies}

The four Delambre formulas are best stated first without division.  This form
remains meaningful even when a traditional quotient presentation would have a
zero denominator.

\begin{theorem}[Delambre--Gauss analogies]
\label{thm:delambre}
For a proper spherical triangle,
\begin{align}
 \sin\frac{A+B}{2}\cos\frac c2
 &=\cos\frac{a-b}{2}\cos\frac C2,                           \tag{7.1}\\
 \sin\frac{A-B}{2}\sin\frac c2
 &=\sin\frac{a-b}{2}\cos\frac C2,                           \tag{7.2}\\
 \cos\frac{A+B}{2}\cos\frac c2
 &=\cos\frac{a+b}{2}\sin\frac C2,                           \tag{7.3}\\
 \cos\frac{A-B}{2}\sin\frac c2
 &=\sin\frac{a+b}{2}\sin\frac C2.                           \tag{7.4}
\end{align}
Cyclic permutation gives twelve cross-multiplied identities.  Whenever the
displayed divisors are apart from zero, these are equivalent to the usual
four quotient analogies.
\end{theorem}

\begin{proof}
We prove (7.1) in detail; the other three use the same calculation with the
corresponding signs.  Squaring and applying the half-angle identities reduces
(7.1) to
\begin{equation}
 (1-\cos(A+B))(1+\cos c)
 =(1+\cos(a-b))(1+\cos C).                                 \tag{7.5}
\end{equation}
The angle cosine law gives
\begin{align*}
 1-\cos(A+B)
 &=1-\cos A\cos B+\sin A\sin B\\
 &=1+\cos C+\sin A\sin B(1-\cos c).
\end{align*}
The side cosine law gives, dually,
\begin{align*}
 1+\cos(a-b)
 &=1+\cos a\cos b+\sin a\sin b\\
 &=1+\cos c+\sin a\sin b(1-\cos C).
\end{align*}
After multiplication by \(1+\cos c\) and \(1+\cos C\), respectively, the
constant terms agree.  The remaining terms agree because
\begin{equation}
 \sin A\sin B\sin^2c=\sin a\sin b\sin^2C,                 \tag{7.6}
\end{equation}
which is the spherical sine law twice.  This proves equality of the squares.
Both sides of (7.1) are nonnegative: \(\cos(c/2),\cos(C/2)>0\),
\(\cos((a-b)/2)>0\) by the strict triangle inequalities, and
\(\sin((A+B)/2)>0\).  Nonnegative-root uniqueness therefore yields (7.1).

For (7.2)--(7.4), square the two sides, replace each half-square by
\((1\pm\cos)/2\), and use the same two cosine laws plus (7.6).  The required
signs follow from the proper-triangle inequalities; signed versions may also
be retained directly when no order simplification is desired.  Every squared
calculation is polynomial after clearing the positive factors.
\end{proof}

The quotient forms most commonly printed are
\begin{align}
 \frac{\sin((A+B)/2)}{\cos(C/2)}
 &=\frac{\cos((a-b)/2)}{\cos(c/2)},                         \tag{7.7}\\
 \frac{\sin((A-B)/2)}{\cos(C/2)}
 &=\frac{\sin((a-b)/2)}{\sin(c/2)},                         \tag{7.8}\\
 \frac{\cos((A+B)/2)}{\sin(C/2)}
 &=\frac{\cos((a+b)/2)}{\cos(c/2)},                         \tag{7.9}\\
 \frac{\cos((A-B)/2)}{\sin(C/2)}
 &=\frac{\sin((a+b)/2)}{\sin(c/2)}.                         \tag{7.10}
\end{align}
Here the half-side and half-angle sines/cosines are positive except for the
possible factors involving \(A-B\) or \(a-b\); those occur only in numerators.

\subsection{Napier's analogies}

\begin{theorem}[Napier analogies]
\label{thm:napier}
With the usual apartness guards,
\begin{align}
 \tan\frac{A+B}{2}
 &=\frac{\cos((a-b)/2)}{\cos((a+b)/2)}\cot\frac C2,         \tag{7.11}\\
 \tan\frac{A-B}{2}
 &=\frac{\sin((a-b)/2)}{\sin((a+b)/2)}\cot\frac C2,         \tag{7.12}\\
 \tan\frac{a+b}{2}
 &=\frac{\cos((A-B)/2)}{\cos((A+B)/2)}\tan\frac c2,        \tag{7.13}\\
 \tan\frac{a-b}{2}
 &=\frac{\sin((A-B)/2)}{\sin((A+B)/2)}\tan\frac c2.        \tag{7.14}
\end{align}
The cyclic permutations give twelve formulas.
\end{theorem}

\begin{proof}
Divide (7.1) by (7.3) to obtain (7.11), and divide (7.2) by (7.4) to obtain
(7.12).  Dividing (7.4) by (7.3) gives (7.13); dividing (7.2) by (7.1) gives
(7.14).  Constructively each division is justified by an explicit apartness
witness.  Alternatively, retain the corresponding cross-multiplied identity,
which follows without any quotient.
\end{proof}

\subsection{The spherical law of tangents}

\begin{corollary}[Law of tangents]
\label{cor:tangent-law}
For each cyclic pair,
\begin{equation}
 \frac{\tan((A-B)/2)}{\tan((A+B)/2)}
 =\frac{\tan((a-b)/2)}{\tan((a+b)/2)},                     \tag{7.15}
\end{equation}
with the required apartness assumptions.  Thus also
\begin{align}
 \frac{\tan((B-C)/2)}{\tan((B+C)/2)}
 &=\frac{\tan((b-c)/2)}{\tan((b+c)/2)},                    \tag{7.16}\\
 \frac{\tan((C-A)/2)}{\tan((C+A)/2)}
 &=\frac{\tan((c-a)/2)}{\tan((c+a)/2)}.                    \tag{7.17}
\end{align}
\end{corollary}

\begin{proof}
Divide (7.12) by (7.11), or equivalently divide (7.14) by (7.13).  The common
factor \(\cot(C/2)\), respectively \(\tan(c/2)\), cancels with its positive
apartness witness.  After cross multiplication, (7.15) is again a finite
polynomial consequence of Theorem~\ref{thm:delambre}.
\end{proof}

\begin{remark}[Constructive preference for cross forms]
The quotient identities are convenient for numerical solution of a triangle,
but their logically strongest versions are the cross-multiplied equations.
They remain valid at symmetric configurations such as \(A=B\) or \(a=b\),
where a printed ratio may become \(0/0\).  The formalization therefore proves
cross forms first and derives quotients only under apartness.
\end{remark}

\section{Cotangent, five-part, Cagnoli, right, and quadrantal rules}
\label{sec:cotangent-special}

\subsection{The six cotangent four-part formulas}

The six parts of a spherical triangle occur cyclically as
\(a,C,b,A,c,B\).  Choosing four consecutive parts gives one inner side, one
inner angle, one outer side and one outer angle.  The general mnemonic is
\begin{equation}
 \cos(\text{inner side})\cos(\text{inner angle})
 =\cot(\text{outer side})\sin(\text{inner side})
  -\cot(\text{outer angle})\sin(\text{inner angle}).        \tag{8.1}
\end{equation}
Explicitly,
\begin{align}
 \cos b\cos C&=\cot a\sin b-\cot A\sin C,                 \tag{CT1}\\
 \cos b\cos A&=\cot c\sin b-\cot C\sin A,                 \tag{CT2}\\
 \cos c\cos A&=\cot b\sin c-\cot B\sin A,                 \tag{CT3}\\
 \cos c\cos B&=\cot a\sin c-\cot A\sin B,                 \tag{CT4}\\
 \cos a\cos B&=\cot c\sin a-\cot C\sin B,                 \tag{CT5}\\
 \cos a\cos C&=\cot b\sin a-\cot B\sin C.                 \tag{CT6}
\end{align}
For a proper triangle every side and angle sine is positive, so all cotangents
have canonical apartness witnesses.

\begin{proof}[Proof of CT1 and the family]
Begin with the cosine law for \(a\), substitute the cosine law for \(c\), and
use the sine law to replace \(\sin c\):
\begin{align*}
 \cos a
 &=\cos b\cos c+\sin b\sin c\cos A\\
 &=\cos b(\cos a\cos b+\sin a\sin b\cos C)
   +\sin b\sin C\sin a\cot A.
\end{align*}
Move the \(\cos a\cos^2b\) term to the left:
\[
 \cos a\sin^2b
 =\cos b\sin a\sin b\cos C
  +\sin b\sin C\sin a\cot A.
\]
Cancel the positive product \(\sin a\sin b\) to obtain CT1.  The same
substitution pattern gives CT3 and CT5.  Applying those three identities to
the polar triangle gives CT2, CT4 and CT6.  Before cancellation, each proof is
an integer polynomial identity, so it is valid without any division.
\end{proof}

Two frequently printed consequences are already visible here:
\begin{align}
 \cos A&=\csc b\,\csc c\,(\cos a-\cos b\cos c),           \tag{8.2}\\
 \cos a\cos C&=\sin a\cot b-\sin C\cot B.                 \tag{8.3}
\end{align}

\subsection{Five-part rules}

Substitute the cosine law for \(b\) into the law for \(a\):
\begin{align*}
 \cos a
 &=(\cos a\cos c+\sin a\sin c\cos B)\cos c
   +\sin b\sin c\cos A.
\end{align*}
After moving \(\cos a\cos^2c\) to the left and canceling the positive factor
\(\sin c\), one obtains
\begin{equation}
 \cos a\sin c
 =\sin a\cos c\cos B+\sin b\cos A.                         \tag{8.4}
\end{equation}
Equivalently,
\begin{equation}
 \sin a\cos B=\cos b\sin c-\sin b\cos c\cos A.            \tag{8.5}
\end{equation}
Cyclic permutation gives
\begin{align}
 \cos b\sin a&=\sin b\cos a\cos C+\sin c\cos B,          \tag{8.6}\\
 \cos c\sin b&=\sin c\cos b\cos A+\sin a\cos C.          \tag{8.7}
\end{align}
Polar duality and reversal of the cyclic order produce the remaining standard
five-part variants.  Their division-free versions require no side conditions.

\subsection{Cagnoli's equation}

\begin{theorem}[Cagnoli]
\label{thm:cagnoli}
The six parts of a spherical triangle satisfy
\begin{equation}
 \sin b\sin c+\cos b\cos c\cos A
 =\sin B\sin C-\cos B\cos C\cos a.                         \tag{8.8}
\end{equation}
There are two cyclic companions.
\end{theorem}

\begin{proof}
Multiply the side cosine law for \(a\) by \(\cos A\), replacing
\(\cos^2A\) by \(1-\sin^2A\):
\begin{equation}
 \cos a\cos A
 =\cos b\cos c\cos A+\sin b\sin c
  -\sin b\sin c\sin^2A.                                   \tag{8.9}
\end{equation}
Multiply the angle cosine law for \(A\) by \(\cos a\):
\begin{equation}
 \cos a\cos A
 =-\cos B\cos C\cos a+\sin B\sin C
  -\sin B\sin C\sin^2a.                                  \tag{8.10}
\end{equation}
The sine law gives
\(
 \sin b\sin c\sin^2A=\sin B\sin C\sin^2a
\).
Subtract the common remainder from (8.9)--(8.10), yielding (8.8).
\end{proof}

\subsection{Right spherical triangles}

Assume \(C=\pi/2\).  Then \(\sin C=1\) and \(\cos C=0\).  The core laws reduce
to Napier's ten rules:
\begin{alignat}{3}
 \text{(R1)}\quad&\cos c=\cos a\cos b,
 &\qquad \text{(R6)}\quad&\tan b=\cos A\tan c,             \tag{8.11}\\
 \text{(R2)}\quad&\sin a=\sin A\sin c,
 &\text{(R7)}\quad&\tan a=\cos B\tan c,                  \tag{8.12}\\
 \text{(R3)}\quad&\sin b=\sin B\sin c,
 &\text{(R8)}\quad&\cos A=\sin B\cos a,                  \tag{8.13}\\
 \text{(R4)}\quad&\tan a=\tan A\sin b,
 &\text{(R9)}\quad&\cos B=\sin A\cos b,                  \tag{8.14}\\
 \text{(R5)}\quad&\tan b=\tan B\sin a,
 &\text{(R10)}\quad&\cos c=\cot A\cot B.                 \tag{8.15}
\end{alignat}
The tangent and cotangent equations are guarded by the corresponding cosine or
sine apartness witnesses.  R1 is the side cosine law; R2--R3 are the sine law;
R8--R9 come from the angle cosine laws.  R4--R7 and R10 follow by guarded
division.  Thus the ten-rule mnemonic introduces no new theorem.

Napier's circle summarizes them constructively.  Write the five parts
\[
 a,\quad \frac\pi2-B,\quad b,\quad \frac\pi2-c,\quad \frac\pi2-A
\]
in cyclic order.  For any three selected parts, the sine of the middle part is
the product of the tangents of the two adjacent parts and also the product of
the cosines of the two opposite parts.  Each assertion is one of
(R1)--(R10) after the complement identities.

\subsection{Quadrantal spherical triangles}

Assume instead that \(c=\pi/2\).  Applying the right-triangle rules to the
polar triangle gives
\begin{alignat}{3}
 \text{(Q1)}\quad&\cos C=-\cos A\cos B,
 &\qquad \text{(Q6)}\quad&\tan B=-\cos a\tan C,            \tag{8.16}\\
 \text{(Q2)}\quad&\sin A=\sin a\sin C,
 &\text{(Q7)}\quad&\tan A=-\cos b\tan C,                   \tag{8.17}\\
 \text{(Q3)}\quad&\sin B=\sin b\sin C,
 &\text{(Q8)}\quad&\cos a=\sin b\cos A,                  \tag{8.18}\\
 \text{(Q4)}\quad&\tan A=\tan a\sin B,
 &\text{(Q9)}\quad&\cos b=\sin a\cos B,                  \tag{8.19}\\
 \text{(Q5)}\quad&\tan B=\tan b\sin A,
 &\text{(Q10)}\quad&\cos C=-\cot a\cot b.                \tag{8.20}
\end{alignat}
Again the quotient forms are derived only under apartness.  The minus signs in
Q1, Q6, Q7 and Q10 arise from
\(\cos(\pi-x)=-\cos x\) under polar duality.

\section{Differential variations and the quantitative planar limit}
\label{sec:differential-limit}

\subsection{Constructive differentiable families}

A differentiable real name in this section means a fixed primitive-recursive
procedure for rational approximations to \(f(t)\) and \(f'(t)\), together with
a primitive-recursive modulus \(\delta(n)\) such that every rational
\(0<\abs h\le2^{-\delta(n)}\) satisfies
\begin{equation}
 \abs{\frac{f(t+h)-f(t)}h-f'(t)}\le2^{-n}.                 \tag{9.1}
\end{equation}
The product and chain rules are finite error estimates on bounded intervals.
The ProveIt trigonometric packet proves the derivative names
\((\sin x)'=\cos x\) and \((\cos x)'=-\sin x\) from the finite exponential
series, with explicit tail schedules.  Thus differentiating the cosine laws is
legitimate in \(\ipra\) for every fixed primitive-recursive family.

\begin{theorem}[Differential variation formulas]
\label{thm:differential}
Along a differentiable family of proper spherical triangles,
\begin{align}
 da&=\cos C\,db+\cos B\,dc+\sin b\sin C\,dA,              \tag{9.2}\\
 db&=\cos A\,dc+\cos C\,da+\sin c\sin A\,dB,              \tag{9.3}\\
 dc&=\cos B\,da+\cos A\,db+\sin a\sin B\,dC.              \tag{9.4}
\end{align}
Here each symbol denotes the derivative name at the fixed parameter value.
\end{theorem}

\begin{proof}
Differentiate the cosine law for \(a\):
\begin{align*}
 -\sin a\,da
 &=(-\sin b\cos c+\cos b\sin c\cos A)\,db\\
 &\quad+(-\cos b\sin c+\sin b\cos c\cos A)\,dc
       -\sin b\sin c\sin A\,dA.
\end{align*}
The cyclic five-part identities give
\begin{align*}
 \sin a\cos C&=\sin b\cos c-\cos b\sin c\cos A,\\
 \sin a\cos B&=\cos b\sin c-\sin b\cos c\cos A.
\end{align*}
The sine law gives
\(
 \sin b\sin c\sin A=\sin a\sin b\sin C
\).
Substitution yields
\[
 \sin a\,da=
 \sin a\bigl(\cos C\,db+\cos B\,dc+\sin b\sin C\,dA\bigr).
\]
Since \(\sin a\) has a positive apartness witness, cancellation proves (9.2).
The other formulas are cyclic.  At the finite-precision level the proof is a
sum of product-rule estimates followed by the same polynomial simplification;
no infinitesimal objects are added to the language of arithmetic.
\end{proof}

The equations may be solved for any three unknown differentials once the other
three are supplied and the relevant determinant is positively separated from
zero.  Constructively, matrix inversion is therefore guarded by an explicit
rational lower bound on that determinant.

\subsection{A uniform small-triangle estimate}

The usual statement that the spherical cosine law tends to the planar cosine
law can be made quantitative.

\begin{theorem}[Planar cosine law with an explicit remainder]
\label{thm:planar-limit}
Suppose
\(0\le a,b,c\le\varepsilon\le1\) and the exact spherical cosine law holds.
Then
\begin{equation}
 \abs{c^2-a^2-b^2+2ab\cos C}\le2\varepsilon^4.             \tag{9.5}
\end{equation}
Cyclic variants hold for the other sides.
\end{theorem}

\begin{proof}
The finite Taylor bounds, valid constructively on \([0,1]\), are
\begin{align}
 \abs{t^2-2(1-\cos t)}&\le\frac{t^4}{12},                 \tag{9.6}\\
 0\le t-\sin t&\le\frac{t^3}{6}.                          \tag{9.7}
\end{align}
Put \(F(t)=2(1-\cos t)\).  The spherical cosine law implies
\begin{equation}
 F(c)=F(a)+F(b)-2(1-\cos a)(1-\cos b)-2\sin a\sin b\cos C.
                                                                    \tag{9.8}
\end{equation}
Therefore
\begin{align*}
 &c^2-a^2-b^2+2ab\cos C\\
 &=[c^2-F(c)]-[a^2-F(a)]-[b^2-F(b)]\\
 &\quad-2(1-\cos a)(1-\cos b)
       +2(ab-\sin a\sin b)\cos C.
\end{align*}
The first line of errors has absolute value at most
\(3\varepsilon^4/12\).  Since \(1-\cos t\le t^2/2\), the next term is at most
\(\varepsilon^4/2\).  Finally,
\[
 \abs{ab-\sin a\sin b}
 \le a\abs{b-\sin b}+\sin b\abs{a-\sin a}
 \le\frac{\varepsilon^4}{3},
\]
so the last term is at most \(2\varepsilon^4/3\).  Their sum is
\(17\varepsilon^4/12<2\varepsilon^4\), proving (9.5).
All Taylor estimates come from finite alternating sums with primitive-recursive
remainder indices.
\end{proof}

For a sphere of radius \(R\), let \(\ell_a=Ra\), etc., and assume
\(\ell_a,\ell_b,\ell_c\le L\le R\).  Multiplying (9.5) by \(R^2\) gives
\begin{equation}
 \abs{\ell_c^2-\ell_a^2-\ell_b^2+2\ell_a\ell_b\cos C}
 \le\frac{2L^4}{R^2}.                                      \tag{9.9}
\end{equation}
Thus the curvature correction is explicitly quadratic in \(1/R\).

\begin{corollary}[Planar sine-law cross relation]
Under the same small-side hypothesis,
\begin{equation}
 \abs{a\sin B-b\sin A}\le\frac{\varepsilon^3}{3},          \tag{9.10}
\end{equation}
with cyclic companions.
\end{corollary}

\begin{proof}
The exact spherical sine law says
\(\sin a\sin B=\sin b\sin A\).  Hence
\[
 a\sin B-b\sin A
 =(a-\sin a)\sin B-(b-\sin b)\sin A,
\]
and (9.7) gives the bound.
\end{proof}

\section{Area, spherical excess, solid angle, and polygons}
\label{sec:area}

\subsection{Constructive spherical area}

A statement about area requires more than the algebraic triangle certificate,
but it still admits a primitive-recursive finite-precision interpretation.
Use the equal-area coordinates
\begin{equation}
 \Phi(\lambda,z)=
 \bigl(\sqrt{1-z^2}\cos\lambda,
       \sqrt{1-z^2}\sin\lambda,z\bigr),
 \quad -\pi\le\lambda\le\pi,\quad -1\le z\le1.            \tag{10.1}
\end{equation}
Away from the two pole lines,
\begin{equation}
 \norm{\partial_\lambda\Phi\crossp\partial_z\Phi}=1.       \tag{10.2}
\end{equation}
Thus surface area is ordinary rectangle area \(d\lambda\,dz\) in this chart.
All coordinate functions have the constructive names from Section~\ref{sec:ipra}.

For a polygon bounded by \(m\) great-circle arcs, subdivide the
\((\lambda,z)\)-rectangle into a dyadic grid.  A cell whose interval evaluation
has a fixed sign against every boundary normal is decidably inside or outside.
The undecided cells lie in a union of strips around the boundary great circles.
The strip of angular half-width \(\delta\) around one great circle has area
\begin{equation}
 4\pi\sin\delta\le4\pi\delta.                              \tag{10.3}
\end{equation}
Therefore the gap between lower and upper cell sums is at most
\(4\pi m\delta\), plus the explicitly bounded coordinate-roundoff margin.
Choosing a dyadic \(\delta\) from the requested output precision gives a
primitive-recursive modulus.  This constructs a rapid area name.  Cellwise
counting proves finite additivity; boundary strips vanish with the same
modulus.  No appeal to a nonconstructive measure-extension theorem is needed.

There is an even simpler computation for a lune.  Choose the polar axis of its
two bounding great circles as the \(z\)-axis.  A lune of angle \(\alpha\)
corresponds to the rectangle
\(0\le\lambda\le\alpha\), \(-1\le z\le1\).  Equations
(10.1)--(10.2) give
\begin{equation}
 \Area(\text{lune of angle }\alpha)=2\alpha.               \tag{10.4}
\end{equation}
At finite precision the Riemann sum is exact because the area density is the
constant \(1\).

\subsection{Girard's theorem}

\begin{theorem}[Girard's spherical excess formula]
\label{thm:girard}
For a proper spherical triangle on the unit sphere,
\begin{equation}
 \Area(ABC)=E:=A+B+C-\pi.                                  \tag{10.5}
\end{equation}
On a sphere of radius \(R\), the area is \(R^2E\).
\end{theorem}

\begin{proof}
The three side great circles partition the sphere into eight spherical
triangles, paired antipodally.  Let the area of the original triangle be
\(T\), and choose one triangle from each of the other three antipodal pairs,
with areas \(X,Y,Z\), so that these four regions form a hemisphere.  Finite
additivity gives
\begin{equation}
 T+X+Y+Z=2\pi.                                              \tag{10.6}
\end{equation}
The interior-angle lune at \(A\) consists of the original triangle together
with one of the regions \(X,Y,Z\); after relabeling take its area to be
\(T+X\).  Similarly the \(B\)- and \(C\)-lunes have areas \(T+Y\) and
\(T+Z\).  By (10.4), their total area is
\begin{equation}
 2A+2B+2C=3T+X+Y+Z.                                        \tag{10.7}
\end{equation}
Use (10.6) to replace \(X+Y+Z\) by \(2\pi-T\).  Then
\[
 2(A+B+C)=2T+2\pi,
\]
so \(T=A+B+C-\pi\).  The proof is a finite dissection identity plus the exact
lune calculation.  Both are valid for the rapid area names just constructed.
Scaling a unit-sphere patch by \(R\) multiplies every area by \(R^2\).
\end{proof}

\subsection{An algebraic solid-angle formula}

The area can also be recovered directly from the vector certificate.  Put
\begin{equation}
 D=1+x+y+z,
 \qquad
 Q=4\cos\frac a2\cos\frac b2\cos\frac c2
   =\sqrt{2(1+x)(1+y)(1+z)}.                               \tag{10.8}
\end{equation}
The polynomial identity
\begin{equation}
 D^2+\tau^2=2(1+x)(1+y)(1+z)=Q^2                          \tag{10.9}
\end{equation}
uses only the Gram relation.

\begin{theorem}[Oriented solid-angle formula]
\label{thm:solid-angle}
For the positive orientation,
\begin{align}
 \cos\frac E2&=\frac{D}{Q},                                \tag{10.10}\\
 \sin\frac E2&=\frac{\tau}{Q},                             \tag{10.11}\\
 E&=2\atanTwo(\tau,D).                                     \tag{10.12}
\end{align}
If \(D\apart0\), then
\begin{equation}
 \tan\frac E2=\frac{\tau}{1+x+y+z}.                       \tag{10.13}
\end{equation}
\end{theorem}

\begin{proof}
Since \(E/2=P-\pi/2\), the addition formulas give
\(\cos(E/2)=\sin P\) and \(\sin(E/2)=-\cos P\).  Expand
\(\sin P\) and \(\cos P\) by successive half-angle addition.  Substitute the
half-angle formulas of Section~\ref{sec:half-formulas}, clear the positive
factor \(4\cos(a/2)\cos(b/2)\cos(c/2)\), and use (4.16).  The two numerators
reduce respectively to \(D\) and \(\tau\).  This is a finite polynomial
normalization after the certified square-root equations are inserted.
Equation (10.9) verifies that the displayed pair is a unit pair.  Since
\(0<E<2\pi\) and \(\tau>0\), the constructive two-chart \(\atanTwo\) operation
selects exactly \(E/2\), proving (10.12).  Guarded division yields (10.13).
\end{proof}

Formula (10.12) is often the most stable computational form of Girard's
excess.  It avoids subtracting nearly equal angles for a small triangle.

\subsection{L'Huilier's formula}

\begin{corollary}[L'Huilier]
\label{cor:lhuillier}
Let \(p=(a+b+c)/2\).  Then
\begin{equation}
 \tan^2\frac E4=
 \tan\frac p2
 \tan\frac{p-a}{2}
 \tan\frac{p-b}{2}
 \tan\frac{p-c}{2},                                       \tag{10.14}
\end{equation}
and the positive square root gives \(\tan(E/4)\).
\end{corollary}

\begin{proof}
From (10.10),
\begin{equation}
 \tan^2\frac E4=\frac{1-\cos(E/2)}{1+\cos(E/2)}
 =\frac{Q-D}{Q+D}.                                         \tag{10.15}
\end{equation}
Write each tangent on the right of (10.14) as a sine divided by a cosine.
Apply product-to-sum twice to the four numerator factors and twice to the four
denominator factors.  After inserting
\(Q=4\cos(a/2)\cos(b/2)\cos(c/2)\), the resulting numerator is
\(Q-D\) and the denominator is \(Q+D\).  Equivalently, cross multiplication
reduces to (10.9) and the three identities
\(2\sin r\sin s=\cos(r-s)-\cos(r+s)\).  All factors in (10.14) are positive
for a proper triangle, so nonnegative-root uniqueness fixes the sign.
\end{proof}

\subsection{Spherical polygons}

\begin{theorem}[Polygonal excess]
\label{thm:polygon-excess}
Let a positively oriented simple geodesic \(n\)-gon on the unit sphere have
interior angles \(\alpha_1,\ldots,\alpha_n\), and suppose a certified
noncrossing triangulation lies in the polygon.  Then
\begin{equation}
 \Area(P)=\sum_{j=1}^n\alpha_j-(n-2)\pi.                   \tag{10.16}
\end{equation}
On a radius-\(R\) sphere, multiply by \(R^2\).
\end{theorem}

\begin{proof}
A triangulation has \(n-2\) triangles.  Finite additivity says the polygonal
area is the sum of their areas.  The triangle angles partition the polygon's
interior angles, including at vertices incident to several diagonals.  Applying
Girard to each triangle therefore gives
\[
 \Area(P)
 =\sum_{j=1}^n\alpha_j-(n-2)\pi.
\]
The entire argument is finite.  The triangulation and its noncrossing/orientation
signs are part of the arithmetic certificate, so no choice of a diagonal is
made inside \(\ipra\).
\end{proof}

\section{Tetrahedral trigonometry}
\label{sec:tetrahedron}

Let \(P_1,P_2,P_3,P_4\in\R^3\) be the vertices of a nondegenerate Euclidean
tetrahedron.  Let \(F_i\) be the face opposite \(P_i\), let \(\Delta_i>0\) be
its area, and let \(\theta_{ij}\in(0,\pi)\) be the interior dihedral angle
along the edge \(P_iP_j\).

\subsection{Oriented face-area vectors}

For each face choose the outward unit normal \(n_i\) and define its doubled
area vector
\begin{equation}
 N_i=2\Delta_i n_i.                                         \tag{11.1}
\end{equation}

\begin{lemma}[Closure of face-area vectors]
\label{lem:face-closure}
The four outward face-area vectors satisfy
\begin{equation}
 N_1+N_2+N_3+N_4=0.                                        \tag{11.2}
\end{equation}
\end{lemma}

\begin{proof}
Translate \(P_4\) to the origin and put
\(a=P_1-P_4\), \(b=P_2-P_4\), \(c=P_3-P_4\), with
\(\detm[a\ b\ c]>0\).  The inward doubled normals of the first three faces
are \(b\crossp c\), \(c\crossp a\), \(a\crossp b\).  The oriented normal of
the face \(P_1P_2P_3\) is
\begin{align*}
 (b-a)\crossp(c-a)
 &=b\crossp c+c\crossp a+a\crossp b.
\end{align*}
It points away from \(P_4\), so its inward normal is the negative of this sum.
Thus the four inward vectors sum to zero; negating all four proves (11.2) for
outward vectors.  This is a coordinate polynomial identity.
\end{proof}

Faces \(F_j\) and \(F_k\) meet along the edge joining the two remaining
vertices \(P_i,P_l\).  The angle between their outward normals is
\(\pi-\theta_{il}\).  Hence
\begin{equation}
 N_j\dotp N_k
 =4\Delta_j\Delta_k\cos(\pi-\theta_{il})
 =-4\Delta_j\Delta_k\cos\theta_{il}.                       \tag{11.3}
\end{equation}
This relation may itself be taken as the constructive definition of the
circle-valued dihedral angle; the positive sine comes from the normalized
cross product of the two face normals.

\subsection{The tetrahedral cosine law}

\begin{theorem}[Cosine-law analogue for a tetrahedron]
\label{thm:tetra-cosine}
For every permutation \(\{i,j,k,l\}=\{1,2,3,4\}\),
\begin{align}
 \Delta_i^2={}&\Delta_j^2+\Delta_k^2+\Delta_l^2\\
 &-2\bigl(
 \Delta_j\Delta_k\cos\theta_{il}
 +\Delta_j\Delta_l\cos\theta_{ik}
 +\Delta_k\Delta_l\cos\theta_{ij}
 \bigr).                                                    \tag{11.4}
\end{align}
\end{theorem}

\begin{proof}
By closure,
\(N_i=-(N_j+N_k+N_l)\).  Take squared norms:
\begin{align*}
 4\Delta_i^2
 &=4\Delta_j^2+4\Delta_k^2+4\Delta_l^2\\
 &\quad+2N_j\dotp N_k+2N_j\dotp N_l+2N_k\dotp N_l.
\end{align*}
Insert (11.3).  The pair \((j,k)\) leaves the edge \((i,l)\), the pair
\((j,l)\) leaves \((i,k)\), and the pair \((k,l)\) leaves \((i,j)\).
Division by \(4\) yields (11.4).
\end{proof}

This proof is the exact three-dimensional analogue of expanding
\(\norm{x+y}^2\) in the planar cosine law.  The ``side lengths'' are now face
areas and the relevant angles are supplements of the angles between outward
normals.

\begin{corollary}[Linear normal relation]
\label{cor:tetra-linear}
Dotting (11.2) with \(n_i\) gives the stronger first-degree relation
\begin{equation}
 \Delta_i=
 \Delta_j\cos\theta_{kl}+
 \Delta_k\cos\theta_{jl}+
 \Delta_l\cos\theta_{jk}.                                 \tag{11.5}
\end{equation}
\end{corollary}

For a regular tetrahedron, all four areas are equal and (11.5) gives
\(3\cos\theta=1\), the familiar \(\cos\theta=1/3\).

\begin{corollary}[Dihedral Gram determinant]
Let \(G\) be the \(4\times4\) Gram matrix of the outward unit normals:
\begin{equation}
 G_{ii}=1,
 \qquad
 G_{ij}=-\cos\theta_{kl}
 \quad(\{i,j,k,l\}=\{1,2,3,4\}).                           \tag{11.6}
\end{equation}
Then
\begin{equation}
 \detm G=0,
 \qquad
 G(\Delta_1,\Delta_2,\Delta_3,\Delta_4)^T=0.               \tag{11.7}
\end{equation}
\end{corollary}

\begin{proof}
Four vectors in \(\R^3\) are linearly dependent, so their Gram determinant
vanishes.  Equation (11.2), divided by \(2\), is precisely the displayed
weighted null relation.
\end{proof}

\subsection{The tetrahedral sine-product identity}

Let the vertices be denoted \(O,A,B,C\).  All angles in the next theorem are
ordinary planar angles inside the triangular faces.

\begin{theorem}[Tetrahedral sine law]
\label{thm:tetra-sine}
\begin{equation}
 \sin\angle OAB\;\sin\angle OBC\;\sin\angle OCA
 =\sin\angle OAC\;\sin\angle OCB\;\sin\angle OBA.          \tag{11.8}
\end{equation}
Choosing each of the four vertices in turn as \(O\) gives four relations, of
which at most three are independent.
\end{theorem}

\begin{proof}
Apply the ordinary planar sine law in the three faces through \(O\):
\begin{align*}
 \frac{\sin\angle OAB}{\sin\angle OBA}&=\frac{OB}{OA},\\
 \frac{\sin\angle OBC}{\sin\angle OCB}&=\frac{OC}{OB},\\
 \frac{\sin\angle OCA}{\sin\angle OAC}&=\frac{OA}{OC}.
\end{align*}
Multiplying the three equations makes all edge lengths cancel, leaving (11.8).
The planar sine laws can be imported from ProveIt's constructive triangle
package; every denominator is positive in a nondegenerate face.
\end{proof}

Multiplying any three of the four vertex-based equations and canceling the
common positive sine factors yields the fourth.  This explains constructively
why only three constraints are independent.

\subsection{The \(\ipra\) interpretation}

The tetrahedral proof is even more elementary than the spherical one.  Encode
the four vertex coordinates by fixed rapid names.  Cross products give the
unnormalized face-area vectors exactly.  A bounded square-root routine gives
\(2\Delta_i=\norm{N_i}\), with positive lower bounds supplied by
nondegeneracy.  Define \(\cos\theta_{ij}\) by the normalized negative dot
product of the two outward normals.  Then (11.2) is a coordinate polynomial
identity and (11.4) is one ring expansion.  Actual angle names are optional;
the cosine law needs only their cosine values.  Consequently the object-level
proof requires no analytic trigonometry beyond normalization by positive
square roots.

\section{Lean formalization plan and the object-language endpoint}
\label{sec:formalization}

The archive contains three Lean files.  They are deliberately separated by
logical role.

All three files are under the archive's \path{Lean/} directory.
\begin{description}[style=nextline,leftmargin=0pt,labelindent=0pt]
\item[\code{SphericalKernel.lean}]
Coordinate-free finite algebra: vectors, dot/cross products, Gram and Lagrange
identities, the core spherical numerators, polar involution, face-area closure,
and the tetrahedral norm expansion.  It is intended to use only Mathlib's ring
tactics.
\item[\code{SphericalTrig.lean}]
An abstract constructive trigonometric system, a record of spherical core data,
and theorem interfaces for cosine, sine, haversine, half-angle, Delambre,
Napier, special-triangle, differential, area, and tetrahedral statements.
\item[\code{IPRABridge.lean}]
The concrete representation boundary for ProveIt's \code{ConciseIPRA}: codecs,
rapid names, apartness, fixed PR symbols, checker formulas, and the final
\code{Prf PRAAx} targets.
\end{description}

No Lean executable is installed in the delivery environment, so these files
were not compiled.  They contain an algebraic prototype and precise theorem
interfaces; they do not constitute a completed object-level \(\ipra\) proof.
The remaining tasks are explicit below.

\subsection{Stage 1: finite vector algebra}

Use a three-coordinate structure over an arbitrary commutative ring:
\begin{lstlisting}[language={}]
structure V3 (R : Type*) where
  x : R; y : R; z : R

def dot   (u v : V3 R) : R := ...
def cross (u v : V3 R) : V3 R := ...
def det3  (u v w : V3 R) : R := dot u (cross v w)
\end{lstlisting}
Prove by \code{ring}:
\begin{enumerate}
\item Lagrange:
\[
 (u\crossp v)\dotp(u\crossp w)
 =(u\dotp u)(v\dotp w)-(u\dotp v)(u\dotp w).
\]
\item Gram:
\[
 \detm(u,v,w)^2=\detm((u_i\dotp u_j)_{i,j=1}^3).
\]
\item The tangent numerator identities and
\((x-yz)^2+\tau^2=(1-y^2)(1-z^2)\).
\item Polar involution:
\((w\crossp u)\crossp(u\crossp v)=\tau u\).
\item The coordinate face-area closure and the expansion of its squared norm.
\end{enumerate}
These theorems require neither an order nor a field.  They are the most stable
part of the formalization and should be compiled first.  Run
\code{#print axioms} on the final declarations; no custom axioms should occur.

\subsection{Stage 2: an algebraic spherical core record}

To avoid square roots in the main identity proofs, define a record over a
commutative ring:
\begin{lstlisting}[language={}]
structure SphericalCore (R : Type*) [CommRing R] where
  x y z sa sb sc tau : R
  cA sA cB sB cC sC : R
  sa_sq : sa^2 + x^2 = 1
  sb_sq : sb^2 + y^2 = 1
  sc_sq : sc^2 + z^2 = 1
  gram  : tau^2 = 1 + 2*x*y*z - x^2 - y^2 - z^2
  cA_num : sb*sc*cA = x - y*z
  cB_num : sc*sa*cB = y - z*x
  cC_num : sa*sb*cC = z - x*y
  sA_num : sb*sc*sA = tau
  sB_num : sc*sa*sB = tau
  sC_num : sa*sb*sC = tau
\end{lstlisting}
The side cosine laws are just \code{cA_num} rearranged.  The angle cosine laws,
unit-pair equations, sine cross-products, five-part rules and Cagnoli identities
are proved by \code{ring} or \code{linear_combination} from these fields.  This
record is also the semantic target of the exact Python checker.

A second record adds order information only where needed:
\begin{lstlisting}[language={}]
structure ProperCore (R : Type*) [LinearOrderedField R]
    extends SphericalCore R where
  sa_pos : 0 < sa; sb_pos : 0 < sb; sc_pos : 0 < sc
  tau_pos : 0 < tau
\end{lstlisting}
It supports guarded cancellation and positive-root uniqueness.  The finite
algebra should remain in the weaker record whenever possible.

\subsection{Stage 3: abstract constructive trigonometry}

Reuse or import the class \code{TrigSystem} from the ProveIt trigonometric
packet.  Its data are an additive angle group, a commutative coefficient ring,
functions \code{sin}, \code{cos}, a quarter turn, the addition laws,
Pythagorean identity and quarter-turn values.  It already supports parity,
subtraction, periods, double-angle, product-to-sum and square half-angle laws.

Define a named spherical triangle with six angle terms and a core realization:
\begin{lstlisting}[language={}]
structure NamedSphericalTriangle (Angle R : Type*)
    [AddCommGroup Angle] [LinearOrderedField R]
    [TrigSystem Angle R] where
  a b c A B C : Angle
  core : ProperCore R
  cos_a : cos a = core.x
  sin_a : sin a = core.sa
  -- cyclic side equations
  cos_A : cos A = core.cA
  sin_A : sin A = core.sA
  -- cyclic angle equations
  half_side_positive : ...
  half_angle_positive : ...
\end{lstlisting}
Then formalize the formula atlas in dependency order:
\begin{enumerate}
\item core cosine/sine laws and polar substitutions;
\item half-square identities, followed by positive-root forms;
\item cross-multiplied Delambre identities;
\item Napier and tangent laws under denominator hypotheses;
\item cotangent, five-part, Cagnoli, right and quadrantal rules;
\item differential identities in a separate module parameterized by a
constructive derivative interface.
\end{enumerate}
The recommended theorem statements expose cross forms first.  Quotient theorems
should accept explicit hypotheses such as \code{cos x != 0} only at the ambient
field level; the object-level version must instead accept an apartness witness.

\subsection{Stage 4: codecs and primitive-recursive operations}

Pin the implementation to the inspected ProveIt revision
\code{a3c053cc6243}.  Within
\code{LeanProofs.ConciseIPRA}, construct actual \code{PR} syntax trees for:
\begin{enumerate}
\item pairing, projections and fixed-length tuples;
\item signed integers, positive-denominator rationals and dyadics;
\item addition, multiplication, powers of two, absolute value and comparison;
\item bounded sums, products, maximum, scan and finite-list folds;
\item three-vector dot, cross and determinant operations;
\item the bounded square-root scan (2.6);
\item finite exponential and arctangent sums with their explicit cutoffs.
\end{enumerate}
Derive the defining equations with \code{PRA.proj}, \code{PRA.comp},
\code{PRA.precZero} and \code{PRA.precSucc}.  Invalid rational codes should
decode to a fixed harmless default so that every symbol remains total.

\subsection{Stage 5: rapid-name arithmetic}

Represent a fixed real input by particular lower- and upper-endpoint
\code{PR 1} terms, an integer magnitude bound, and derivations of the interval
regularity formulas.  Do not make a Lean function \code{Nat -> Rat} the object
being quantified over.  For each operation, compile a new particular PR term
and prove a finite rational error estimate at a free precision index.

The essential numerical lemmas are:
\begin{align*}
 |xy-x'y'|&\le |x||y-y'|+|y'||x-x'|,\\
 |x^{-1}-y^{-1}|&\le 2|x-y|/d^2\quad(|x|,|y|\ge d/2),\\
 |\sqrt x-\sqrt y|&\le\sqrt{|x-y|},\\
 \left|\sum_{j=N}^{M}\frac{z^j}{j!}\right|&\le2^{-p}
 \quad\text{for the prescribed }N(B,p).
\end{align*}
Every induction invariant is a Boolean combination of natural equalities and
inequalities, encoded by characteristic functions, so it is quantifier free.
Only after the fixed-precision theorem is obtained should universal
introduction quantify over the precision.

\subsection{Stage 6: an object-language polynomial certificate checker}

Encode sparse multivariate polynomials as finite lists of integer coefficients
and exponent vectors.  Implement normalization by insertion, merging like
monomials and deleting zero coefficients.  Prove inside \(\PRAAx\):
\begin{enumerate}
\item evaluation commutes with polynomial addition and multiplication;
\item normalization preserves evaluation;
\item a normalized zero polynomial evaluates to zero;
\item each stored reduction certificate from the core ideal is sound.
\end{enumerate}
Prove these by structural primitive recursion and quantifier-free induction on
list lengths.  The external Python and Lean ring checks are specifications and
regressions.  Acceptance requires the internal soundness derivation, not merely
a Boolean checker returning \code{true}.

The existing \code{EqKernel}, \code{EqArith} and normalization modules in
\code{ConciseIPRA} should be reused where their interfaces match; otherwise the
small specialized normalizer above is sufficient.

\subsection{Stage 7: geometry-to-core realization}

Given fixed coordinate names for \(u,v,w\):
\begin{enumerate}
\item compute dot/cross/determinant names;
\item prove the unit-vector and Gram identities from the polynomial checker;
\item compute positive side-sine names with bounded square root;
\item use supplied lower bounds to construct reciprocals;
\item compute \((c_A,s_A)\), etc., and prove the unit-pair equations;
\item construct the angle names once using the two-chart \(\atanTwo\) routine;
\item prove their sine/cosine realization and branch intervals.
\end{enumerate}
The chart choice is made at one fixed certified precision and retained.  A
fresh undecidable branch choice at every output precision would not define a
Cauchy name.

\subsection{Stage 8: area and finite polygons}

Represent a geodesic polygon by a finite list of oriented great-circle normals
and a triangulation certificate.  Implement the lower/upper dyadic cell counts
in the equal-area chart (10.1).  The strip estimate (10.3) supplies the modulus.
Prove finite additivity at the cell level, then derive the exact lune formula.
The eight-region partition underlying Girard is finite combinatorics; encode
its incidence table directly.  Polygonal excess follows by primitive recursion
on the triangle list.

As a computationally lighter alternative for triangles, define the area name
by the solid-angle formula (10.12) and prove its agreement with the cell-area
name through Girard.  The solid-angle route is preferable for executable
geometry; the lune route establishes the intended surface-area semantics.

\subsection{Stage 9: tetrahedra}

The tetrahedral module needs only coordinate algebra and positive square
roots.  Construct the four oriented doubled face vectors, prove their sum is
zero by \code{ring}, define each dihedral cosine by a normalized negative dot
product, and expand the squared closure.  The sine-product identity imports the
already formalized planar sine law on three faces.  No spherical angle names or
area integration are required.

\subsection{Stage 10: the final derivability theorem}

For each fixed collection of input PR terms, assemble a closed formula
\(\mathsf{SphHyp}\to\mathsf{SphConclusion}\).  The final target has the shape
\begin{lstlisting}[language={}]
theorem spherical_atlas_derivable :
  Prf PRAAx [] (Fm.imp sphericalHypotheses sphericalConclusions)
\end{lstlisting}
where every auxiliary function symbol is a concrete compiled \code{PR} term.
After this theorem exists, apply \code{pra_theorem_iff} or
\code{pra_iff_native} to transfer it to the desired concise native
presentation.  The transport theorem is the final step, not a substitute for
the derivation.

\subsection{Validation commands}

In a checkout with Lean/Mathlib 4.32.0 and the pinned ProveIt revision, the
initial validation sequence is intended to be
\begin{lstlisting}[language=bash]
lake env lean Lean/SphericalKernel.lean
lake env lean Lean/SphericalTrig.lean
lake build ConciseIPRA.NativeBiInterpretation
lake env lean Lean/IPRABridge.lean
\end{lstlisting}
Resolve elaboration or tactic-version differences before treating the files as
checked.  Then audit the final ambient kernel theorems with \code{#print axioms}
and, separately, require an actual closed \code{Prf PRAAx} term for the
arithmetized theorem.  These are distinct acceptance criteria.

\section{Coverage, constructive solution procedures, and limitations}
\label{sec:coverage}

\subsection{Coverage of the requested source pages}

The following table records the mathematical coverage.  ``Cyclic'' means that
the displayed generator is proved once and all cyclic relabelings are obtained
by the same proof term or a finite renaming transformation.

\begin{longtable}{@{}p{0.24\textwidth}p{0.17\textwidth}p{0.49\textwidth}@{}}
\toprule
Family & Status & Constructive content \\
\midrule
\endfirsthead
\toprule
Family & Status & Constructive content \\
\midrule
\endhead
Spherical side cosine law & proved, 3 cyclic & Vector tangent projection;
polynomial numerator form and guarded quotient form.\\
Spherical sine law & proved & Scalar triple product; common ratio and all
cross-multiplied forms.\\
Supplemental/angle cosine law & proved, 3 cyclic & Direct Gram reduction and
polar-dual proof.\\
Polar triangle & proved & Explicit normalized cross products, branch signs,
and involution.\\
Half-angle formulas & proved, 9 cyclic instances & Square identities first;
positive-root forms from branch certificates.\\
Half-side formulas & proved, 9 cyclic instances & Polar duals of half-angle
formulas, including the requested formula.\\
Half-tangent forms & proved, 6 cyclic instances & Rational formulas in the six
half-tangents, with denominator guards.\\
Haversines & proved, cyclic and dual & Stable navigation form and inverse
form.\\
Delambre--Gauss analogies & proved, 12 cyclic instances & Cross forms first;
quotients under apartness.\\
Napier analogies & proved, 12 cyclic instances & Guarded quotients of
Delambre identities.\\
Spherical tangent law & proved, 3 cyclic instances & Quotients of Napier
analogies; cross form valid at symmetric cases.\\
Cotangent/four-part rules & proved, 6 & Cosine-law substitution and polar
duality.\\
Five-part rules & proved, cyclic/dual family & Finite substitution and
cancellation; includes the MathWorld miscellaneous identities.\\
Cagnoli equation & proved, 3 cyclic & Side/angle cosine laws plus the sine
law.\\
Right-triangle rules & proved, 10 & Direct specialization at \(C=\pi/2\).\\
Quadrantal rules & proved, 10 & Polar duals of right-triangle rules.\\
Differential variations & proved, 3 & Constructive derivatives with explicit
moduli; no infinitesimals as objects.\\
Planar limit & proved quantitatively & Explicit error
\(\le2\varepsilon^4\) for the cosine law and
\(\le\varepsilon^3/3\) for the sine cross relation.\\
Girard and polygon excess & proved & Equal-area chart, finite lune dissection,
explicit boundary-strip modulus.\\
Solid-angle and L'Huilier formulas & proved & Two-chart \(\atanTwo\), Gram
factorization, positive branches.\\
Tetrahedral cosine law & proved, 4 vertex choices & Closure of oriented
face-area vectors and a norm expansion.\\
Tetrahedral sine product & proved, 4 vertex choices & Product of three planar
face sine laws.\\
\bottomrule
\end{longtable}

Counting cyclic instances separately, the appendices contain well over one
hundred explicit identities.  The exact regression artifact verifies 68 finite generator identities in 14
families, with zero failures.  It does not attempt to certify every quotient
and square-root presentation independently; instead it checks the polynomial
generators from which those presentations are derived.  The branch,
apartness, convergence and object-language obligations remain explicit parts
of the mathematical and formalization arguments.

\subsection{Constructive solution of spherical triangles}

The formula atlas also yields finite algorithms for the classical data cases.
An ``inverse'' operation below means a two-chart constructive angle name plus a
branch interval; it never means an unguarded classical choice.

\begin{enumerate}[label=\textup{Case \arabic*.}]
\item \textbf{SSS.}  Given \(a,b,c\) with certified proper-triangle
inequalities, use the positive half-angle formulas (6.4)--(6.6) to obtain
\(A,B,C\) uniquely in \((0,\pi)\).  This avoids inverse-cosine ambiguity.
\item \textbf{SAS.}  Given \(b,c,A\), compute \(\cos a\) from (4.1), compute
its positive sine from the Gram/radical identity, then construct
\(a=\atanTwo(\sin a,\cos a)\).  Reduce to SSS.
\item \textbf{SSA.}  Given \(a,b,A\), the sine law determines \(\sin B\).
There can be one or two branch-certified values \(B\in(0,\pi)\).  Enumerate the
two candidates \(B\) and \(\pi-B\), and retain exactly those passing a cosine
law and the proper-triangle inequalities.  The search is finite and decidable
at the certificate level.
\item \textbf{ASA.}  Given \(A,C,b\), use two cotangent four-part equations to
obtain the remaining sides, with determinant apartness, then apply the sine
law for the last angle.
\item \textbf{AAS.}  This is the polar/dual version of SSA and has the same
one-or-two-solution branch ledger.
\item \textbf{AAA.}  Use the positive half-side formulas (6.9)--(6.13) to
obtain \(a,b,c\) uniquely.
\item \textbf{SSAA.}  Given two sides and their opposite angles, Napier's
analogies compute the half-sums and half-differences of the remaining pair,
again under explicit denominator and branch certificates.
\end{enumerate}

Each algorithm consists of bounded rational operations, bounded square-root
search and finitely many chart/branch tests.  Therefore each fixed certified
instance has a primitive-recursive real-name implementation.  The ambiguous
SSA/AAS cases return a finite list rather than invoking an excluded-middle
choice of one solution.

\subsection{What is not claimed}

The report does not claim that the accompanying Lean files are compiled, nor
that a completed \code{Prf PRAAx} term has already been constructed.  It also
does not claim a universal internal solver quantified over arbitrary
real-number functions or arbitrary PR program codes.  The result is the
first-order schema described in Section~\ref{sec:ipra}, with a concrete route
to internal derivations.

The area construction is developed far enough to give a primitive-recursive
modulus for certified finite geodesic polygons.  A general constructive
surface-measure library for arbitrary subsets of the sphere is neither needed
nor asserted.  Similarly, degenerate triangles, antipodal vertices and
unoriented polygons require separate signed or limiting conventions; the
proper/nondegenerate hypotheses are maintained explicitly instead of being
silently erased by classical case splits.

\appendix
\section{Compact formula atlas}
\label{app:atlas}

This appendix collects the formulas in one place.  Every quotient carries the
apartness conditions explained in the body, and every square root is the
positive branch justified by properness.  Put
\[
 p=\frac{a+b+c}{2},\qquad P=\frac{A+B+C}{2},\qquad
 E=A+B+C-\pi.
\]

\subsection{Primary laws}

\begin{align*}
\cos a&=\cos b\cos c+\sin b\sin c\cos A,\\
\cos b&=\cos c\cos a+\sin c\sin a\cos B,\\
\cos c&=\cos a\cos b+\sin a\sin b\cos C;\\[0.3em]
\cos A&=-\cos B\cos C+\sin B\sin C\cos a,\\
\cos B&=-\cos C\cos A+\sin C\sin A\cos b,\\
\cos C&=-\cos A\cos B+\sin A\sin B\cos c;\\[0.3em]
\frac{\sin A}{\sin a}&=
\frac{\sin B}{\sin b}=
\frac{\sin C}{\sin c}.
\end{align*}

Rearrangements:
\begin{align*}
\cos A&=\frac{\cos a-\cos b\cos c}{\sin b\sin c},&
\cos a&=\frac{\cos A+\cos B\cos C}{\sin B\sin C},
\end{align*}
plus cyclic versions.

\subsection{Polar triangle}

\begin{align*}
a^*&=\pi-A,&b^*&=\pi-B,&c^*&=\pi-C,\\
A^*&=\pi-a,&B^*&=\pi-b,&C^*&=\pi-c,
\end{align*}
and polarity is involutive.

\subsection{Half-angle and half-side families}

For \((A,a)\), with cyclic companions:
\begin{align*}
\sin\frac A2&=\sqrt{\frac{\sin(p-b)\sin(p-c)}{\sin b\sin c}},\\
\cos\frac A2&=\sqrt{\frac{\sin p\sin(p-a)}{\sin b\sin c}},\\
\tan\frac A2&=\sqrt{\frac{\sin(p-b)\sin(p-c)}{\sin p\sin(p-a)}};\\[0.4em]
\sin\frac a2&=\sqrt{-\frac{\cos P\cos(P-A)}{\sin B\sin C}},\\
\cos\frac a2&=\sqrt{\frac{\cos(P-B)\cos(P-C)}{\sin B\sin C}},\\
\tan\frac a2&=\sqrt{-\frac{\cos P\cos(P-A)}{\cos(P-B)\cos(P-C)}}.
\end{align*}

Common factors:
\begin{align*}
k&=[\sin p\sin(p-a)\sin(p-b)\sin(p-c)]^{-1/2},\\
\tan\frac A2&=k\sin(p-b)\sin(p-c),\\
K&=\sqrt{-\frac{\cos P}{\cos(P-A)\cos(P-B)\cos(P-C)}},\\
\tan\frac a2&=K\cos(P-A),
\end{align*}
plus cyclic companions.

Direct vector forms, where
\(x=\cos a,y=\cos b,z=\cos c\) and
\(\tau=\detm[u\ v\ w]\):
\begin{align*}
\tan\frac A2&=\frac{\tau}{\sin b\sin c+\cos a-\cos b\cos c},\\
\tan\frac a2&=\frac{\sin a}{1+\cos a}.
\end{align*}

With \(t_a=\tan(a/2)\) and similarly for all six parts:
\begin{align*}
t_A^2&=
 \frac{(t_a-t_b+t_c+t_at_bt_c)(t_a+t_b-t_c+t_at_bt_c)}
 {(t_a+t_b+t_c-t_at_bt_c)(-t_a+t_b+t_c+t_at_bt_c)},\\
t_a^2&=
 \frac{(t_Bt_C+t_Ct_A+t_At_B-1)
       (-t_Bt_C+t_Ct_A+t_At_B+1)}
 {(t_Bt_C-t_Ct_A+t_At_B+1)
       (t_Bt_C+t_Ct_A-t_At_B+1)}.
\end{align*}

\subsection{Haversines}

\begin{align*}
\hav a&=\hav(b-c)+\sin b\sin c\,\hav A,\\
\hav b&=\hav(c-a)+\sin c\sin a\,\hav B,\\
\hav c&=\hav(a-b)+\sin a\sin b\,\hav C;\\
\hav A&=\hav(\pi-B-C)+\sin B\sin C\,\hav a,
\end{align*}
with the two cyclic duals and guarded inverse forms.

\subsection{Delambre--Gauss analogies}

\begin{align*}
\frac{\sin((A+B)/2)}{\cos(C/2)}&=
 \frac{\cos((a-b)/2)}{\cos(c/2)},\\
\frac{\sin((A-B)/2)}{\cos(C/2)}&=
 \frac{\sin((a-b)/2)}{\sin(c/2)},\\
\frac{\cos((A+B)/2)}{\sin(C/2)}&=
 \frac{\cos((a+b)/2)}{\cos(c/2)},\\
\frac{\cos((A-B)/2)}{\sin(C/2)}&=
 \frac{\sin((a+b)/2)}{\sin(c/2)},
\end{align*}
plus eight cyclic companions.  The cross-multiplied forms are the primary
constructive statements.

\subsection{Napier analogies and tangent law}

\begin{align*}
\tan\frac{A+B}{2}&=
 \frac{\cos((a-b)/2)}{\cos((a+b)/2)}\cot\frac C2,\\
\tan\frac{A-B}{2}&=
 \frac{\sin((a-b)/2)}{\sin((a+b)/2)}\cot\frac C2,\\
\tan\frac{a+b}{2}&=
 \frac{\cos((A-B)/2)}{\cos((A+B)/2)}\tan\frac c2,\\
\tan\frac{a-b}{2}&=
 \frac{\sin((A-B)/2)}{\sin((A+B)/2)}\tan\frac c2,
\end{align*}
plus cyclic companions, and
\[
\frac{\tan((A-B)/2)}{\tan((A+B)/2)}
 =\frac{\tan((a-b)/2)}{\tan((a+b)/2)}
\]
plus its two cyclic companions.

\subsection{Cotangent and five-part formulas}

\begin{align*}
\cos b\cos C&=\cot a\sin b-\cot A\sin C,\\
\cos b\cos A&=\cot c\sin b-\cot C\sin A,\\
\cos c\cos A&=\cot b\sin c-\cot B\sin A,\\
\cos c\cos B&=\cot a\sin c-\cot A\sin B,\\
\cos a\cos B&=\cot c\sin a-\cot C\sin B,\\
\cos a\cos C&=\cot b\sin a-\cot B\sin C.
\end{align*}
Representative five-part forms:
\begin{align*}
\cos a\sin c&=\sin a\cos c\cos B+\sin b\cos A,\\
\sin a\cos B&=\cos b\sin c-\sin b\cos c\cos A,
\end{align*}
with cyclic and polar companions.

Cagnoli:
\[
\sin b\sin c+\cos b\cos c\cos A
 =\sin B\sin C-\cos B\cos C\cos a,
\]
plus two cyclic companions.

\subsection{Right triangle, \(C=\pi/2\)}

\begin{align*}
\cos c&=\cos a\cos b,&
\sin a&=\sin A\sin c,&
\sin b&=\sin B\sin c,\\
\tan a&=\tan A\sin b,&
\tan b&=\tan B\sin a,&
\tan b&=\cos A\tan c,\\
\tan a&=\cos B\tan c,&
\cos A&=\sin B\cos a,&
\cos B&=\sin A\cos b,\\
\cos c&=\cot A\cot B.&&
\end{align*}

\subsection{Quadrantal triangle, \(c=\pi/2\)}

\begin{align*}
\cos C&=-\cos A\cos B,&
\sin A&=\sin a\sin C,&
\sin B&=\sin b\sin C,\\
\tan A&=\tan a\sin B,&
\tan B&=\tan b\sin A,&
\tan B&=-\cos a\tan C,\\
\tan A&=-\cos b\tan C,&
\cos a&=\sin b\cos A,&
\cos b&=\sin a\cos B,\\
\cos C&=-\cot a\cot b.&&
\end{align*}

\subsection{Differential, area, and solid-angle formulas}

\begin{align*}
da&=\cos C\,db+\cos B\,dc+\sin b\sin C\,dA,\\
db&=\cos A\,dc+\cos C\,da+\sin c\sin A\,dB,\\
dc&=\cos B\,da+\cos A\,db+\sin a\sin B\,dC.\\[0.3em]
\Area(ABC)&=E=A+B+C-\pi,\\
\Area(P_n)&=\sum_{j=1}^n\alpha_j-(n-2)\pi,\\
E&=2\atanTwo\bigl(\tau,1+\cos a+\cos b+\cos c\bigr),\\
\tan^2\frac E4&=
 \tan\frac p2\tan\frac{p-a}{2}\tan\frac{p-b}{2}\tan\frac{p-c}{2}.
\end{align*}

\subsection{Tetrahedral formulas}

For \(\{i,j,k,l\}=\{1,2,3,4\}\):
\begin{align*}
\Delta_i^2={}&\Delta_j^2+\Delta_k^2+\Delta_l^2\\
&-2\bigl(\Delta_j\Delta_k\cos\theta_{il}
+\Delta_j\Delta_l\cos\theta_{ik}
+\Delta_k\Delta_l\cos\theta_{ij}\bigr),\\
\Delta_i={}&\Delta_j\cos\theta_{kl}
+\Delta_k\cos\theta_{jl}
+\Delta_l\cos\theta_{jk}.
\end{align*}
For vertices \(O,A,B,C\):
\[
\sin\angle OAB\sin\angle OBC\sin\angle OCA
=\sin\angle OAC\sin\angle OCB\sin\angle OBA.
\]

\section{Representative internal \(\ipra\) targets}
\label{app:ipra-target}

This appendix spells out the distinction between the finite exact kernel and a
statement about represented real coordinates.

\subsection{A single exact rational checker theorem}

Let \(w\in\N\) encode a tuple of signed rational values for all symbols in
(4.16).  Define primitive-recursive characteristic functions:
\begin{itemize}
\item \(\mathsf{CoreCert}(w)=1\) when all denominator codes are positive and
all equations in (4.16) hold after cross multiplication;
\item \(\mathsf{CoreConsequences}(w)=1\) when the cross-multiplied side and
angle cosine laws, sine law, unit-angle equations, five-part identities and
Cagnoli equations all hold.
\end{itemize}
Both functions use only bounded tuple decoding and natural arithmetic.  The
first exact target is
\begin{equation}
 \ipra\vdash
 \forall w\,
 \bigl(\mathsf{CoreCert}(w)=1\to
       \mathsf{CoreConsequences}(w)=1\bigr).                \tag{B.1}
\end{equation}
A verified polynomial normalizer can produce the derivation of (B.1).  No real
numbers or limits occur here.

For the tetrahedron, an analogous checker
\(\mathsf{TetCert}\) verifies four coordinate face vectors, their closure,
positive squared norms and the normalized dihedral dot equations.  Its target
is
\begin{equation}
 \ipra\vdash
 \forall w\,
 \bigl(\mathsf{TetCert}(w)=1\to
       \mathsf{TetCosine}(w)=1\bigr).                       \tag{B.2}
\end{equation}

\subsection{A fixed-name finite-precision theorem}

Now fix, in the metatheory, particular primitive-recursive programs for the
nine coordinates of \(u,v,w\), their moduli and bounds.  The programs are not
object-language variables.  Compile particular PR symbols for all output names
and let \(\mathsf{ApproxHyp}(n,\vec q)\) assert that the rational tuple
\(\vec q\) consists of the finitely many input approximations requested by the
precision schedule for output precision \(n\).  Let
\(\mathsf{CosineError}(n,\vec q)\) be the natural-number inequality obtained
by clearing positive denominators in
\begin{equation}
 \abs{\cos a-
 (\cos b\cos c+\sin b\sin c\cos A)}\le2^{-n}.              \tag{B.3}
\end{equation}
The actual internal theorem has the form
\begin{equation}
 \ipra\vdash
 \forall n\,\forall\vec q\,
 \bigl(\mathsf{ApproxHyp}(n,\vec q)
 \to\mathsf{CosineError}(n,\vec q)\bigr).                  \tag{B.4}
\end{equation}
The proof at free \(n\) instantiates input regularity only at explicitly
computed indices, invokes the finite polynomial identity, and closes a rational
error estimate.  Its induction calls are on bounded loops and use
quantifier-free invariants.  Universal introduction of \(n\) comes last.

The sine law, half formulas and all remaining identities use the same shape.
A quotient conclusion additionally assumes a numerical apartness certificate;
a positive-root conclusion includes the bounded-square-root bracket and the
branch interval.

\subsection{Lean object-language shape}

Using the actual repository syntax, the intended endpoint is schematically
\begin{lstlisting}[language={}]
open LeanProofs.ConciseIPRA

def sphericalHypotheses : Fm 0 := ...
def sphericalConclusions : Fm 0 := ...

theorem sphericalAtlas :
  Prf PRAAx [] (Fm.imp sphericalHypotheses sphericalConclusions) := by
  -- concrete PR defining equations
  -- qf induction for bounded arithmetic and error loops
  -- internally sound polynomial certificates
  -- ordinary intuitionistic introduction/elimination rules
  ...
\end{lstlisting}
The ellipsis denotes work, not an assumed theorem.  The supplied
\code{IPRABridge.lean} defines the data interfaces and target signatures so
that the remaining implementation obligations are visible rather than hidden
behind host-level axioms.

\subsection{Semantic presentation theorem}

After the internal theorem is complete, the existing ProveIt
bi-interpretation can present it in the concise native language.  Separately,
a semantic theorem states that standard rapid dyadic names satisfying the
premises denote the ordinary geometric quantities and hence validate the
familiar real-number formulas.  The direction of justification is therefore
\[
 \text{finite arithmetic derivation}
 \longrightarrow
 \text{name equality at every precision}
 \longrightarrow
 \text{ordinary geometric presentation},
\]
not the reverse.


\begin{thebibliography}{99}
\bibitem{wiki-slc}
Wikipedia contributors, ``Spherical law of cosines,'' accessed 15 November 2026.
\url{https://en.wikipedia.org/wiki/Spherical_law_of_cosines}

\bibitem{wiki-sines}
Wikipedia contributors, ``Law of sines: Spherical law of sines,'' accessed 15 November 2026.
\url{https://en.wikipedia.org/wiki/Law_of_sines#Spherical_law_of_sines}

\bibitem{wiki-halfside}
Wikipedia contributors, ``Half-side formula,'' accessed 15 November 2026.
\url{https://en.wikipedia.org/wiki/Half-side_formula}

\bibitem{wiki-costet}
Wikipedia contributors, ``Law of cosines: Analogue for tetrahedra,'' accessed 15 November 2026.
\url{https://en.wikipedia.org/wiki/Law_of_cosines#Analogue_for_tetrahedra}

\bibitem{wiki-sphtrig}
Wikipedia contributors, ``Spherical trigonometry,'' accessed 15 November 2026.
\url{https://en.wikipedia.org/wiki/Spherical_trigonometry}

\bibitem{mathworld}
E. W. Weisstein, ``Spherical Trigonometry,'' MathWorld--A Wolfram Web Resource,
accessed 15 November 2026.
\url{https://mathworld.wolfram.com/SphericalTrigonometry.html}

\bibitem{todhunter}
I. Todhunter, \emph{Spherical Trigonometry}, 5th ed., Macmillan, 1886.

\bibitem{bishopbridges}
E. Bishop and D. Bridges, \emph{Constructive Analysis}, Springer, 1985.

\bibitem{bridgesrichman}
D. Bridges and F. Richman, \emph{Varieties of Constructive Mathematics},
Cambridge University Press, 1987.

\bibitem{proveit}
V. Reshetnikov, \emph{ProveIt}, repository snapshot
\code{a3c053cc6243}, 15 November 2026.
\url{https://github.com/VladimirReshetnikov/ProveIt}

\bibitem{trigpacket}
ProveIt, \emph{A Constructive Atlas of Trigonometric Identities in Intuitionistic
Primitive Recursive Arithmetic}, directory
\path{Logic/PrimitiveRecursiveArithmetic/docs/proofs/analysis/trigonometric_identities_ipra_2026-11-15}, 2026.

\bibitem{trianglepacket}
ProveIt, \emph{Heron, sine, and cosine laws in intuitionistic PRA}, directory
\path{Logic/PrimitiveRecursiveArithmetic/docs/proofs/geometry/triangle_formulas_ipra}, 2026.

\bibitem{rodriguespacket}
ProveIt, \emph{Euler's rotation theorem and Rodrigues' formula in iPRA}, directory
\path{Logic/PrimitiveRecursiveArithmetic/docs/proofs/geometry/euler_rodrigues_ipra}, 2026.
\end{thebibliography}

\end{document}
